\documentclass[../main.tex]{subfiles}
\begin{document}

\chapter{GPUアーキテクチャの最適化 II}
\lessoninfo{
  video={M7，レッスン 1--3（GPU Virtual Memory；GPU Virtual Memory - UVA；GPU Warp Scheduling）},
  sources={Power ら \cite{power2014x86}（GPU MMU）；Ausavarungnirun ら \cite{ausavarungnirun2017mosaic}（複数のページサイズ）；Li ら \cite{li2021tlb}（マルチ GPU の TLB）；Ganguly ら \cite{ganguly2019interplay}（ユニファイドメモリにおけるプリフェッチと追い出し）；Kim ら \cite{kim2020batch}（フォールトのバッチ処理）；Rogers ら \cite{rogers2012ccws}（キャッシュ考慮型ウェーブフロントスケジューリング）；Narasiman ら \cite{narasiman2011twolevel} および Gebhart ら \cite{gebhart2011energy}（2レベルスケジューリング）；Jog ら \cite{jog2013orchestrated}（プリフェッチを考慮したスケジューリング），\cite{jog2013owl}（OWL）},
  textbook={A18 \cite{aamodt2018gpu} ch.~3，ch.~4；K12 \cite{kim2012gpgpu} §1.4.1 GPU Pipeline；CPG \cite{nvidia2024cudaguide} §4.2 Hardware Multithreading，§19 Unified Memory Programming；BPG \cite{nvidia2026cudabestpractices} §10.1.4 Unified Virtual Addressing},
  keywords={仮想メモリ，アドレス変換，TLB，多段ページテーブル，ページテーブルウォーカ，ラージページ，セグメントマッピング，コピー後実行，共有仮想メモリ，オンデマンドページング，ファーフォールト，ツリーベース近傍プリフェッチ，オーバーサブスクリプション，ワープスケジューラ，ラウンドロビン，GTO，2レベルスケジューリング，ワーキングセット，キャッシュ考慮型ウェーブフロントスケジューリング},
}

第6章では SM の内部を変えた．ワープがダイバージェント分岐をどう
乗り越えるか，そしてそのレジスタをどう格納するかである．本章では
SM の仕組みのさらに2つの部分を扱う．最初の2つの節はメモリアドレスに
関するものである．CPU とデータ構造を共有する GPU には仮想メモリが
必要であり，ワープ全体のアドレスを一度に変換することは，CPU コアでは
決して起こらない形で変換ハードウェアに負荷をかける．CPU と GPU が
1つのアドレス空間を共有すると，ページの割り当てと移動も，カーネルの
実行中に PCIe を越えて，要求に応じて行わなければならない．第3の節は
第4章のパイプラインのスケジューラ段に戻り，複数のワープが準備完了の
ときにどのワープを発行すべきかを問う．この選択は，メモリレイテンシが
どれだけうまく隠蔽されるか，そして各ワープが L1 キャッシュをどれだけ
保持できるかを決める．第6章と同様に，本章は研究論文を軸に構成されて
おり，各節はその拠り所となる論文を挙げる．

\section{GPU の仮想メモリ}
\label{sec:l07-vm}

仮想メモリには2つの側面がある．プログラムが発行するすべてのアドレスの
変換と，その背後にある物理ページの割り当てである．\source{M7，GPU 仮想
メモリ；A18 ch.~4；\cite{power2014x86}；CPG §19.2.2.1 Memory Paging
and Page Sizes．}CPU ではオペレーティングシステムがページを割り当てる．
GPU では，割り当てはもともと CPU が GPU ドライバとともに，
\texttt{cudaMalloc}，\texttt{cudaMallocManaged}，\texttt{cudaMemcpy}
を通じて行っていた．これには\cref{sec:l07-uva}で立ち返る．本節では
アドレス変換を復習し，スレッドではなくワープがアドレスを発行するときに
何が異なるかを示し，変換を安価にする3つの方法を紹介する．

\subsection{アドレス変換}

\begin{definition}[仮想メモリ]
  \label{def:l07-vm}
  \term[かそうめもり]{仮想メモリ}[virtual memory]のもとでは，
  プログラムは私的なアドレス空間の中で
  \term[かそうあどれす]{仮想アドレス}[virtual address]を発行し，
  ハードウェアがアクセスのたびにそれをメモリ上の
  \term[ぶつりあどれす]{物理アドレス}[physical address]に対応づける．
  この対応づけを
  \term[あどれすへんかん]{アドレス変換}[address translation]という．
  対応は\term[ぺーじ]{ページ}[page]の粒度で保持される．ページとは
  アドレス空間の固定サイズで整列したブロックであり，典型的には
  \SI{4}{\kibi\byte} または \SI{2}{\mebi\byte} である．
\end{definition}

仮想アドレスは2つのフィールドに分かれる（\cref{fig:l07-translation}）．
上位ビットが
\term[かそうぺーじばんごう]{仮想ページ番号}[virtual page number (VPN)]（VPN），
下位ビットがページ内の位置を表す
\term[ぺーじおふせっと]{ページオフセット}[page offset]である．
変換は VPN を
\term[ぶつりふれーむばんごう]{物理フレーム番号}[physical frame number (PFN)]（PFN）
で置き換え，オフセットは変えない．PFN は物理ページ番号（PPN）とも呼ばれる．
$2^k$ バイトのページは $k$ ビットのオフセットを持つ．\index{PFN|see{物理フレーム番号}}\index{VPN|see{仮想ページ番号}}\index{PPN|see{物理フレーム番号}}

\begin{marginfigure}
  % Address translation (Lesson 7): a virtual address splits into a
% virtual page number and a page offset; translation replaces the VPN
% by a physical frame number and keeps the offset unchanged.  Field
% widths are those of 4 KiB pages with a 48-bit virtual address.
% Redraws the outline figure m7l1_translation.
% Usage: % Address translation (Lesson 7): a virtual address splits into a
% virtual page number and a page offset; translation replaces the VPN
% by a physical frame number and keeps the offset unchanged.  Field
% widths are those of 4 KiB pages with a 48-bit virtual address.
% Redraws the outline figure m7l1_translation.
% Usage: % Address translation (Lesson 7): a virtual address splits into a
% virtual page number and a page offset; translation replaces the VPN
% by a physical frame number and keeps the offset unchanged.  Field
% widths are those of 4 KiB pages with a 48-bit virtual address.
% Redraws the outline figure m7l1_translation.
% Usage: \input{figures/l07-translation}  (in a marginfigure)
% Width: margin (50 mm); the drawing is about 46 mm wide.
\begin{tikzpicture}[x=1mm, y=1mm, font=\scriptsize\sffamily,
    >={Stealth[length=1.6mm]}]
  % virtual address
  \node[anchor=east] at (-0.5,23) {VA};
  \draw[fill=ln-accent!20] (0,20) rectangle (28,26) node[midway] {VPN};
  \draw[fill=ln-box] (28,20) rectangle (44,26) node[midway] {\iflangja{オフセット}{offset}};
  \node[anchor=south, text=ln-muted] at (0.8,26) {47};
  \node[anchor=south, text=ln-muted] at (28,26) {12\,11};
  \node[anchor=south, text=ln-muted] at (43.2,26) {0};
  % translation
  \draw[->] (14,20) -- (14,13) node[midway, right, align=left, text=ln-muted]
    {TLB /\\\iflangja{ページテーブル}{page table}};
  \draw[->] (36,20) -- (36,13) node[midway, right, text=ln-muted]
    {\iflangja{不変}{copied}};
  % physical address
  \node[anchor=east] at (-0.5,10) {PA};
  \draw[fill=ln-tip!25] (0,7) rectangle (28,13) node[midway] {PFN};
  \draw[fill=ln-box] (28,7) rectangle (44,13) node[midway] {\iflangja{オフセット}{offset}};
\end{tikzpicture}
  (in a marginfigure)
% Width: margin (50 mm); the drawing is about 46 mm wide.
\begin{tikzpicture}[x=1mm, y=1mm, font=\scriptsize\sffamily,
    >={Stealth[length=1.6mm]}]
  % virtual address
  \node[anchor=east] at (-0.5,23) {VA};
  \draw[fill=ln-accent!20] (0,20) rectangle (28,26) node[midway] {VPN};
  \draw[fill=ln-box] (28,20) rectangle (44,26) node[midway] {\iflangja{オフセット}{offset}};
  \node[anchor=south, text=ln-muted] at (0.8,26) {47};
  \node[anchor=south, text=ln-muted] at (28,26) {12\,11};
  \node[anchor=south, text=ln-muted] at (43.2,26) {0};
  % translation
  \draw[->] (14,20) -- (14,13) node[midway, right, align=left, text=ln-muted]
    {TLB /\\\iflangja{ページテーブル}{page table}};
  \draw[->] (36,20) -- (36,13) node[midway, right, text=ln-muted]
    {\iflangja{不変}{copied}};
  % physical address
  \node[anchor=east] at (-0.5,10) {PA};
  \draw[fill=ln-tip!25] (0,7) rectangle (28,13) node[midway] {PFN};
  \draw[fill=ln-box] (28,7) rectangle (44,13) node[midway] {\iflangja{オフセット}{offset}};
\end{tikzpicture}
  (in a marginfigure)
% Width: margin (50 mm); the drawing is about 46 mm wide.
\begin{tikzpicture}[x=1mm, y=1mm, font=\scriptsize\sffamily,
    >={Stealth[length=1.6mm]}]
  % virtual address
  \node[anchor=east] at (-0.5,23) {VA};
  \draw[fill=ln-accent!20] (0,20) rectangle (28,26) node[midway] {VPN};
  \draw[fill=ln-box] (28,20) rectangle (44,26) node[midway] {\iflangja{オフセット}{offset}};
  \node[anchor=south, text=ln-muted] at (0.8,26) {47};
  \node[anchor=south, text=ln-muted] at (28,26) {12\,11};
  \node[anchor=south, text=ln-muted] at (43.2,26) {0};
  % translation
  \draw[->] (14,20) -- (14,13) node[midway, right, align=left, text=ln-muted]
    {TLB /\\\iflangja{ページテーブル}{page table}};
  \draw[->] (36,20) -- (36,13) node[midway, right, text=ln-muted]
    {\iflangja{不変}{copied}};
  % physical address
  \node[anchor=east] at (-0.5,10) {PA};
  \draw[fill=ln-tip!25] (0,7) rectangle (28,13) node[midway] {PFN};
  \draw[fill=ln-box] (28,7) rectangle (44,13) node[midway] {\iflangja{オフセット}{offset}};
\end{tikzpicture}

  \caption{変換はページオフセットを保ち，VPN を PFN で置き換える
    （\SI{4}{\kibi\byte} ページと 48 ビットの仮想アドレスの場合の
    ビット位置）．}
  \label{fig:l07-translation}
\end{marginfigure}

\begin{definition}[ページテーブル，TLB，ページフォールト]
  \label{def:l07-tlb}
  \term[ぺーじてーぶる]{ページテーブル}[page table]は，対応を保持する
  メモリ上のデータ構造である．各
  \term[ぺーじてーぶるえんとり]{ページテーブルエントリ}[page table entry (PTE)]\index{PTE|see{ページテーブルエントリ}}（PTE）
  は1つの VPN を PFN に対応づけ，有効ビットや権限ビットなどの
  メタデータを持つ．
  \term[てぃーえるびー]{TLB}[TLB (translation lookaside buffer)]\index{translation lookaside buffer|see{TLB}}
  （translation lookaside buffer）は最近使われた変換の小さなキャッシュ
  である．\term[てぃーえるびーみす]{TLBミス}[TLB miss]は TLB が変換を
  保持していないときに起こり，ページテーブルを読まなければならない．
  \term[ぺーじふぉーると]{ページフォールト}[page fault]は，物理ページが
  まだ割り当てられていないためにページテーブルが有効な対応を保持して
  いないときに起こり，オペレーティングシステムがページを割り当てなければ
  ならない．
\end{definition}

\caution{TLB ミスとページフォールトは別の事象である．TLB ミスは
ページテーブルを読むハードウェア（または短いハンドラ）が解決する．
ページフォールトにはオペレーティングシステムが必要であり，GPU では
ホストが必要である（\cref{sec:l07-uva}）．}%
TLB が存在するのは PTE そのものがメモリ上にあるからである．TLB が
なければ，すべてのロードはまず変換を見つけるために1回以上のメモリ
アクセスを必要とする．

\subsection{多段ページテーブル}

VPN で索引づける単一のテーブルは，$v$ ビットの VPN に対して，
プログラムがそのページを使うかどうかにかかわらず $2^{v}$ 個のエントリを
必要とする．

\begin{definition}[多段ページテーブルとページテーブルウォーク]
  \label{def:l07-multilevel}
  \term[ただんぺーじてーぶる]{多段ページテーブル}[multi-level page table]
  は VPN を段ごとに1つずつ，複数のインデックスに分割する．最上段の
  インデックスは最上段のテーブルのエントリを選び，そのエントリは次の段の
  テーブルのベースアドレスを保持する．そのテーブルのエントリは次の
  インデックスで選ばれ，以下同様に続き，最下段のエントリが PTE である．
  この連鎖をたどることが
  \term[ぺーじてーぶるうぉーく]{ページテーブルウォーク}[page table walk]
  である．これはハードウェアの状態機械である
  \term[ぺーじてーぶるうぉーか]{ページテーブルウォーカ}[page table walker (PTW)]（PTW）
  か，ソフトウェアのミスハンドラが行う．使われるどのアドレスからも
  到達しないテーブルは割り当てられない．\index{PTW|see{ページテーブルウォーカ}}
\end{definition}

\cref{fig:l07-page-walk}は x86-64，ARM，GPU が用いる4段の形を示す．\margin{4段の
x86-64：48 ビットの仮想アドレス，9 ビットのインデックス4つと 12 ビットの
オフセット．各テーブルは \SI{8}{\byte} のエントリ 512 個で，ちょうど
\SI{4}{\kibi\byte} ページ1枚である．ルートレジスタは CR3 である．}
ウォークはルートレジスタから始まり，第4段のインデックスを使って第3段の
テーブルのベースを読み，以下同様に第1段のエントリ，すなわち PTE まで
下る．PFN とオフセットを合わせたものが物理アドレスである．節約した空間の
代価は時間である．TLB でミスした1回の変換には，少なくとも4回の依存
したメモリアクセスがかかる．

\begin{figure}[tb]
  \centering
  % Four-level page table walk (Lesson 7), x86-64 layout: a 48-bit
% virtual address holds four 9-bit indices and a 12-bit page offset.
% CR3 points to the level-4 table; each index selects an entry that
% holds the base of the next table; the level-1 entry is the PTE whose
% PFN, joined with the offset, forms the physical address.
% Redraws the outline figure m7l1_fourleveltlb (bit widths re-derived
% for x86-64, slide erratum of outline Appendix C).
% Usage: % Four-level page table walk (Lesson 7), x86-64 layout: a 48-bit
% virtual address holds four 9-bit indices and a 12-bit page offset.
% CR3 points to the level-4 table; each index selects an entry that
% holds the base of the next table; the level-1 entry is the PTE whose
% PFN, joined with the offset, forms the physical address.
% Redraws the outline figure m7l1_fourleveltlb (bit widths re-derived
% for x86-64, slide erratum of outline Appendix C).
% Usage: % Four-level page table walk (Lesson 7), x86-64 layout: a 48-bit
% virtual address holds four 9-bit indices and a 12-bit page offset.
% CR3 points to the level-4 table; each index selects an entry that
% holds the base of the next table; the level-1 entry is the PTE whose
% PFN, joined with the offset, forms the physical address.
% Redraws the outline figure m7l1_fourleveltlb (bit widths re-derived
% for x86-64, slide erratum of outline Appendix C).
% Usage: \input{figures/l07-page-walk}
% Width: body (120 mm); the drawing is about 116 mm wide.
\begin{tikzpicture}[x=0.95mm, y=1mm, font=\scriptsize\sffamily,
    >={Stealth[length=1.6mm]}]
  % virtual address fields
  \node[anchor=east] at (-1,47) {VA};
  \foreach \i/\lab/\bits in {0/L4/47--39, 1/L3/38--30, 2/L2/29--21, 3/L1/20--12} {
    \draw[fill=ln-accent!20] (\i*22,44) rectangle ++(22,6)
      node[midway] {\lab\ \iflangja{インデックス}{index}};
    \node[anchor=south, text=ln-muted] at (\i*22+11,50) {\bits};
  }
  \draw[fill=ln-box] (88,44) rectangle (112,50)
    node[midway] {\iflangja{オフセット}{offset}};
  \node[anchor=south, text=ln-muted] at (100,50) {11--0};
  % tables: three rows each, the middle one selected
  \foreach \i/\lab/\entry in {0/L4/{L3 base}, 1/L3/{L2 base}, 2/L2/{L1 base}, 3/L1/{PTE: PFN}} {
    \draw[fill=white] (\i*22+3,32) rectangle ++(16,-4) node[midway] {$\cdots$};
    \draw[fill=ln-tip!25] (\i*22+3,28) rectangle ++(16,-5) node[midway] {\entry};
    \draw[fill=white] (\i*22+3,23) rectangle ++(16,-4) node[midway] {$\cdots$};
    \node[text=ln-muted] at (\i*22+11,16.5) {\lab\ \iflangja{テーブル}{table}};
    \draw[->] (\i*22+11,44) -- (\i*22+11,32);
  }
  % next-table pointers
  \foreach \i in {0,1,2} {
    \draw[->, ln-accent, thick] (\i*22+19,25.5) -- ++(2,0) |- (\i*22+25,31);
  }
  % CR3
  \draw[->, ln-accent, thick] (-2,30) node[left, text=black] {CR3} -- (3,30);
  % physical address
  \node[anchor=east] at (61,6) {PA};
  \draw[fill=ln-tip!25] (62,3) rectangle (88,9) node[midway] {PFN};
  \draw[fill=ln-box] (88,3) rectangle (112,9) node[midway] {\iflangja{オフセット}{offset}};
  \draw[->, ln-accent, thick] (85,25.5) -- (87,25.5) -- (87,9);
  \draw[->] (100,44) -- (100,9);
\end{tikzpicture}

% Width: body (120 mm); the drawing is about 116 mm wide.
\begin{tikzpicture}[x=0.95mm, y=1mm, font=\scriptsize\sffamily,
    >={Stealth[length=1.6mm]}]
  % virtual address fields
  \node[anchor=east] at (-1,47) {VA};
  \foreach \i/\lab/\bits in {0/L4/47--39, 1/L3/38--30, 2/L2/29--21, 3/L1/20--12} {
    \draw[fill=ln-accent!20] (\i*22,44) rectangle ++(22,6)
      node[midway] {\lab\ \iflangja{インデックス}{index}};
    \node[anchor=south, text=ln-muted] at (\i*22+11,50) {\bits};
  }
  \draw[fill=ln-box] (88,44) rectangle (112,50)
    node[midway] {\iflangja{オフセット}{offset}};
  \node[anchor=south, text=ln-muted] at (100,50) {11--0};
  % tables: three rows each, the middle one selected
  \foreach \i/\lab/\entry in {0/L4/{L3 base}, 1/L3/{L2 base}, 2/L2/{L1 base}, 3/L1/{PTE: PFN}} {
    \draw[fill=white] (\i*22+3,32) rectangle ++(16,-4) node[midway] {$\cdots$};
    \draw[fill=ln-tip!25] (\i*22+3,28) rectangle ++(16,-5) node[midway] {\entry};
    \draw[fill=white] (\i*22+3,23) rectangle ++(16,-4) node[midway] {$\cdots$};
    \node[text=ln-muted] at (\i*22+11,16.5) {\lab\ \iflangja{テーブル}{table}};
    \draw[->] (\i*22+11,44) -- (\i*22+11,32);
  }
  % next-table pointers
  \foreach \i in {0,1,2} {
    \draw[->, ln-accent, thick] (\i*22+19,25.5) -- ++(2,0) |- (\i*22+25,31);
  }
  % CR3
  \draw[->, ln-accent, thick] (-2,30) node[left, text=black] {CR3} -- (3,30);
  % physical address
  \node[anchor=east] at (61,6) {PA};
  \draw[fill=ln-tip!25] (62,3) rectangle (88,9) node[midway] {PFN};
  \draw[fill=ln-box] (88,3) rectangle (112,9) node[midway] {\iflangja{オフセット}{offset}};
  \draw[->, ln-accent, thick] (85,25.5) -- (87,25.5) -- (87,9);
  \draw[->] (100,44) -- (100,9);
\end{tikzpicture}

% Width: body (120 mm); the drawing is about 116 mm wide.
\begin{tikzpicture}[x=0.95mm, y=1mm, font=\scriptsize\sffamily,
    >={Stealth[length=1.6mm]}]
  % virtual address fields
  \node[anchor=east] at (-1,47) {VA};
  \foreach \i/\lab/\bits in {0/L4/47--39, 1/L3/38--30, 2/L2/29--21, 3/L1/20--12} {
    \draw[fill=ln-accent!20] (\i*22,44) rectangle ++(22,6)
      node[midway] {\lab\ \iflangja{インデックス}{index}};
    \node[anchor=south, text=ln-muted] at (\i*22+11,50) {\bits};
  }
  \draw[fill=ln-box] (88,44) rectangle (112,50)
    node[midway] {\iflangja{オフセット}{offset}};
  \node[anchor=south, text=ln-muted] at (100,50) {11--0};
  % tables: three rows each, the middle one selected
  \foreach \i/\lab/\entry in {0/L4/{L3 base}, 1/L3/{L2 base}, 2/L2/{L1 base}, 3/L1/{PTE: PFN}} {
    \draw[fill=white] (\i*22+3,32) rectangle ++(16,-4) node[midway] {$\cdots$};
    \draw[fill=ln-tip!25] (\i*22+3,28) rectangle ++(16,-5) node[midway] {\entry};
    \draw[fill=white] (\i*22+3,23) rectangle ++(16,-4) node[midway] {$\cdots$};
    \node[text=ln-muted] at (\i*22+11,16.5) {\lab\ \iflangja{テーブル}{table}};
    \draw[->] (\i*22+11,44) -- (\i*22+11,32);
  }
  % next-table pointers
  \foreach \i in {0,1,2} {
    \draw[->, ln-accent, thick] (\i*22+19,25.5) -- ++(2,0) |- (\i*22+25,31);
  }
  % CR3
  \draw[->, ln-accent, thick] (-2,30) node[left, text=black] {CR3} -- (3,30);
  % physical address
  \node[anchor=east] at (61,6) {PA};
  \draw[fill=ln-tip!25] (62,3) rectangle (88,9) node[midway] {PFN};
  \draw[fill=ln-box] (88,3) rectangle (112,9) node[midway] {\iflangja{オフセット}{offset}};
  \draw[->, ln-accent, thick] (85,25.5) -- (87,25.5) -- (87,9);
  \draw[->] (100,44) -- (100,9);
\end{tikzpicture}

  \caption{x86-64 のレイアウトによる4段のページテーブルウォーク．各
    インデックスは次のテーブルのベースを保持するエントリを選ぶ．第1段の
    エントリが PTE であり，その PFN をオフセットと結合したものが物理
    アドレスである．}
  \label{fig:l07-page-walk}
\end{figure}

\begin{example}
  \label{ex:l07-table-size}
  x86-64 のレイアウト，すなわち 36 ビットの VPN と \SI{8}{\byte} の
  エントリを考える．
  \begin{itemize}
    \item \emph{1段}：$2^{36} \approx \num{6.9e10}$ 個のエントリ，
      プロセスあたり \SI{512}{\gibi\byte} のページテーブルであり，
      そのほぼすべてが一度も使われないページのためのものである．
    \item \emph{4段，最良の場合}：すべてのページが1つの第1段テーブルの
      下にあるプログラムは各段に1つのテーブルを必要とし，
      $2^9+2^9+2^9+2^9 = \num{2048}$ 個のエントリが
      \SI{16}{\kibi\byte} に収まり，最大 512 ページ
      （\SI{2}{\mebi\byte}）を対応づける．
    \item \emph{疎なプログラム}：アドレス空間内で遠く離れた2つの配列は
      第4段テーブルの下に2本の連鎖を必要とし，$1+2\cdot3 = 7$ 個の
      テーブル，それでも \SI{28}{\kibi\byte} にすぎない．
  \end{itemize}
  数千のエントリが数百億のエントリに取って代わる．\caution{単一段の
テーブルの大きさは $2^{\text{VPN bits}}$ であって
$2^{\text{address bits}}$ ではない．オフセットのビットはテーブルの
索引に使われない．}
\end{example}

\subsection{GPU におけるアドレス変換}

初期の GPU には仮想メモリがないか，固定の1対1の対応しかなかった．
GPU がプログラマブルになり，CPU とデータを共有し始めると，仮想メモリが
必要になった．最も目に見える例が第5章の統合仮想アドレッシングで
ある．\margin{Power らは，400 レーンを持つ AMD A10 APU は
\enquote{can require up to 400 unique translations in a single cycle}
であると指摘している \cite{power2014x86}．}変換は CPU と同じように
働くが，1つの違いがある．それは\emph{アドレスコアレッシング}の後に
行われることである．

\begin{figure}[tb]
  \centering
  % Address translation in the GPU pipeline (Lesson 7).  Top: the
% pipeline stages of Lesson 4.  Bottom: the memory stage expanded; the
% coalescer merges the warp's requests, the L1 data cache is accessed
% with virtual addresses while the TLB translates, a TLB miss goes to
% the page table walker (or a software miss handler), and the L2 cache
% is accessed with physical addresses.
% Redraws the outline figure m7l1_gpu_vm.
% Usage: % Address translation in the GPU pipeline (Lesson 7).  Top: the
% pipeline stages of Lesson 4.  Bottom: the memory stage expanded; the
% coalescer merges the warp's requests, the L1 data cache is accessed
% with virtual addresses while the TLB translates, a TLB miss goes to
% the page table walker (or a software miss handler), and the L2 cache
% is accessed with physical addresses.
% Redraws the outline figure m7l1_gpu_vm.
% Usage: % Address translation in the GPU pipeline (Lesson 7).  Top: the
% pipeline stages of Lesson 4.  Bottom: the memory stage expanded; the
% coalescer merges the warp's requests, the L1 data cache is accessed
% with virtual addresses while the TLB translates, a TLB miss goes to
% the page table walker (or a software miss handler), and the L2 cache
% is accessed with physical addresses.
% Redraws the outline figure m7l1_gpu_vm.
% Usage: \input{figures/l07-gpu-mmu}
% Width: body (120 mm); the drawing is about 114 mm wide.
\begin{tikzpicture}[x=1mm, y=1mm, font=\scriptsize\sffamily,
    >={Stealth[length=1.6mm]},
    st/.style={draw, fill=ln-box, minimum width=15mm, minimum height=8mm,
               align=center, inner sep=1pt},
    un/.style={draw, minimum width=24mm, minimum height=8mm,
               align=center, inner sep=1pt}]
  % pipeline
  \foreach \lab [count=\i from 0] in
      {\iflangja{フェッチ}{fetch}, \iflangja{デコード}{decode}, RF,
       \iflangja{スケジュール}{schedule}, \iflangja{実行\\アドレス生成}{execute\\addr.\ gen.},
       \iflangja{メモリ}{memory}, \iflangja{ライトバック}{write-back}} {
    \node[st] (s\i) at (\i*16.5+7.5,44) {\lab};
  }
  \node[st, fill=ln-accent!25] at (s5) {\iflangja{メモリ}{memory}};
  \foreach \i [evaluate=\i as \j using int(\i+1)] in {0,...,5} {
    \draw[->] (s\i) -- (s\j);
  }
  % memory stage expanded
  \draw[dashed, ln-muted] (0,36) rectangle (114,0);
  \draw[dashed, ln-muted] (s5.south west) -- (0,36);
  \draw[dashed, ln-muted] (s5.south east) -- (114,36);
  \node[un, fill=ln-accent!20] (co) at (14,18) {\iflangja{アドレス\\コアレッサ}{address\\coalescer}};
  \node[un, fill=ln-tip!20] (l1) at (48,27) {\iflangja{L1 D キャッシュ\\（仮想アドレス）}{L1 D-cache\\(virtual address)}};
  \node[un, fill=ln-source!20] (tlb) at (48,9) {L1 TLB};
  \node[un, fill=ln-caution!15] (ptw) at (95,9) {\iflangja{PTW または\\TLB ミスハンドラ}{PTW or\\TLB miss handler}};
  \node[un, fill=ln-tip!20] (l2) at (95,27) {\iflangja{L2 キャッシュ\\（物理アドレス）}{L2 cache\\(physical address)}};
  \draw[->] (co.east) -- ++(4,0) |- (l1.west);
  \draw[->] (co.east) -- ++(4,0) |- (tlb.west);
  \draw[->] (l1) -- node[above, text=ln-muted] {\iflangja{ミス}{miss}} (l2);
  \draw[->] (tlb) -- node[above, text=ln-muted] {\iflangja{ミス}{miss}} (ptw);
  \draw[->] (tlb.north east) -- node[pos=0.45, above left=-1pt, text=ln-muted] {PA} (l2.south west);
  \draw[->] (ptw) -- node[right, text=ln-muted] {\iflangja{メモリ上の PTE}{PTEs in memory}} (l2);
\end{tikzpicture}

% Width: body (120 mm); the drawing is about 114 mm wide.
\begin{tikzpicture}[x=1mm, y=1mm, font=\scriptsize\sffamily,
    >={Stealth[length=1.6mm]},
    st/.style={draw, fill=ln-box, minimum width=15mm, minimum height=8mm,
               align=center, inner sep=1pt},
    un/.style={draw, minimum width=24mm, minimum height=8mm,
               align=center, inner sep=1pt}]
  % pipeline
  \foreach \lab [count=\i from 0] in
      {\iflangja{フェッチ}{fetch}, \iflangja{デコード}{decode}, RF,
       \iflangja{スケジュール}{schedule}, \iflangja{実行\\アドレス生成}{execute\\addr.\ gen.},
       \iflangja{メモリ}{memory}, \iflangja{ライトバック}{write-back}} {
    \node[st] (s\i) at (\i*16.5+7.5,44) {\lab};
  }
  \node[st, fill=ln-accent!25] at (s5) {\iflangja{メモリ}{memory}};
  \foreach \i [evaluate=\i as \j using int(\i+1)] in {0,...,5} {
    \draw[->] (s\i) -- (s\j);
  }
  % memory stage expanded
  \draw[dashed, ln-muted] (0,36) rectangle (114,0);
  \draw[dashed, ln-muted] (s5.south west) -- (0,36);
  \draw[dashed, ln-muted] (s5.south east) -- (114,36);
  \node[un, fill=ln-accent!20] (co) at (14,18) {\iflangja{アドレス\\コアレッサ}{address\\coalescer}};
  \node[un, fill=ln-tip!20] (l1) at (48,27) {\iflangja{L1 D キャッシュ\\（仮想アドレス）}{L1 D-cache\\(virtual address)}};
  \node[un, fill=ln-source!20] (tlb) at (48,9) {L1 TLB};
  \node[un, fill=ln-caution!15] (ptw) at (95,9) {\iflangja{PTW または\\TLB ミスハンドラ}{PTW or\\TLB miss handler}};
  \node[un, fill=ln-tip!20] (l2) at (95,27) {\iflangja{L2 キャッシュ\\（物理アドレス）}{L2 cache\\(physical address)}};
  \draw[->] (co.east) -- ++(4,0) |- (l1.west);
  \draw[->] (co.east) -- ++(4,0) |- (tlb.west);
  \draw[->] (l1) -- node[above, text=ln-muted] {\iflangja{ミス}{miss}} (l2);
  \draw[->] (tlb) -- node[above, text=ln-muted] {\iflangja{ミス}{miss}} (ptw);
  \draw[->] (tlb.north east) -- node[pos=0.45, above left=-1pt, text=ln-muted] {PA} (l2.south west);
  \draw[->] (ptw) -- node[right, text=ln-muted] {\iflangja{メモリ上の PTE}{PTEs in memory}} (l2);
\end{tikzpicture}

% Width: body (120 mm); the drawing is about 114 mm wide.
\begin{tikzpicture}[x=1mm, y=1mm, font=\scriptsize\sffamily,
    >={Stealth[length=1.6mm]},
    st/.style={draw, fill=ln-box, minimum width=15mm, minimum height=8mm,
               align=center, inner sep=1pt},
    un/.style={draw, minimum width=24mm, minimum height=8mm,
               align=center, inner sep=1pt}]
  % pipeline
  \foreach \lab [count=\i from 0] in
      {\iflangja{フェッチ}{fetch}, \iflangja{デコード}{decode}, RF,
       \iflangja{スケジュール}{schedule}, \iflangja{実行\\アドレス生成}{execute\\addr.\ gen.},
       \iflangja{メモリ}{memory}, \iflangja{ライトバック}{write-back}} {
    \node[st] (s\i) at (\i*16.5+7.5,44) {\lab};
  }
  \node[st, fill=ln-accent!25] at (s5) {\iflangja{メモリ}{memory}};
  \foreach \i [evaluate=\i as \j using int(\i+1)] in {0,...,5} {
    \draw[->] (s\i) -- (s\j);
  }
  % memory stage expanded
  \draw[dashed, ln-muted] (0,36) rectangle (114,0);
  \draw[dashed, ln-muted] (s5.south west) -- (0,36);
  \draw[dashed, ln-muted] (s5.south east) -- (114,36);
  \node[un, fill=ln-accent!20] (co) at (14,18) {\iflangja{アドレス\\コアレッサ}{address\\coalescer}};
  \node[un, fill=ln-tip!20] (l1) at (48,27) {\iflangja{L1 D キャッシュ\\（仮想アドレス）}{L1 D-cache\\(virtual address)}};
  \node[un, fill=ln-source!20] (tlb) at (48,9) {L1 TLB};
  \node[un, fill=ln-caution!15] (ptw) at (95,9) {\iflangja{PTW または\\TLB ミスハンドラ}{PTW or\\TLB miss handler}};
  \node[un, fill=ln-tip!20] (l2) at (95,27) {\iflangja{L2 キャッシュ\\（物理アドレス）}{L2 cache\\(physical address)}};
  \draw[->] (co.east) -- ++(4,0) |- (l1.west);
  \draw[->] (co.east) -- ++(4,0) |- (tlb.west);
  \draw[->] (l1) -- node[above, text=ln-muted] {\iflangja{ミス}{miss}} (l2);
  \draw[->] (tlb) -- node[above, text=ln-muted] {\iflangja{ミス}{miss}} (ptw);
  \draw[->] (tlb.north east) -- node[pos=0.45, above left=-1pt, text=ln-muted] {PA} (l2.south west);
  \draw[->] (ptw) -- node[right, text=ln-muted] {\iflangja{メモリ上の PTE}{PTEs in memory}} (l2);
\end{tikzpicture}

  \caption{アドレス変換を含む GPU パイプラインのメモリ段．コアレッサが
    ワープの要求を併合する．TLB が変換する間，L1 データキャッシュは仮想
    アドレスでアクセスされる．TLB ミスは PTW またはミスハンドラへ送られる．
    L2 キャッシュは物理アドレスでアクセスされる．}
  \label{fig:l07-gpu-mmu}
\end{figure}

第4章のパイプラインでは，メモリ段は今や\cref{fig:l07-gpu-mmu}の
ステップを含む．コアレッサがワープのスレッドの要求を併合し（第4章），
その結果の各要求は TLB で引かれる変換を必要とする．ミスは PTW か TLB
ミスハンドラが解決する．この構成では，L1 データキャッシュは仮想アドレスで，
すべての SM が共有する L2 キャッシュは物理アドレスでアクセスされる．\margin{両方の
配置が存在する．Power らの概念実証の MMU はコアレッサの後，L1 の
\emph{前}で変換するので，その L1 は物理アドレスでアクセスされる
\cite{power2014x86}．}

ワープ命令がいくつの変換を必要とするかはコアレッシングが決める．
同じ命令の要求を含む2つのページには，要求どうしがどれだけ近くても
2回の TLB 参照が必要であり，1つのページ内の2つの要求には1回で足りる．

\begin{example}
  \label{ex:l07-tlb-coalescing}
  8スレッドのワープ，\SI{64}{\byte} のメモリ要求，\SI{256}{\byte} の
  おもちゃのページサイズを考える．VPN はアドレスを 8 ビット右シフト
  したものである．\margin{\SI{256}{\byte} のおもちゃのページは数を
小さく保つためのものである．\SI{4}{\kibi\byte} ページでは，
\SI{4}{\kibi\byte} の行を持つ行優先行列の列のように，キロバイト単位の
ストライドで同じ効果が現れる．}\texttt{R1} を介したロードは連続する
  語 64, 68, \ldots, 92 を読む．1つの要求，1つのページ，1回の TLB
  参照である．\texttt{R2} を介したロードは \SI{160}{\byte} の
  ストライドで読む．
  \begin{center}
    \small
    \begin{tabular}{@{}lcccccccc@{}}
      \toprule
      スレッド & 0 & 1 & 2 & 3 & 4 & 5 & 6 & 7 \\
      \midrule
      アドレス & 0 & 160 & 320 & 480 & 640 & 800 & 960 & 1120 \\
      16進 & \texttt{000} & \texttt{0A0} & \texttt{140} & \texttt{1E0} & \texttt{280} & \texttt{320} & \texttt{3C0} & \texttt{460} \\
      VPN & 0 & 0 & 1 & 1 & 2 & 3 & 3 & 4 \\
      \bottomrule
    \end{tabular}
  \end{center}
  どの2つのアドレスも1つの \SI{64}{\byte} 要求に収まらないので，
  このロードは8つの要求を発行する．それらは5つの異なるページに触れる
  ので，同じページへの要求を併合した後でも，1つの命令に5回の TLB
  参照が必要である．\SI{256}{\byte} 以上のストライドでは，すべての
  スレッドがそれぞれ自分のページに触れる．
\end{example}

最悪の場合，コアレスされないロードはワープのスレッド数と同じだけの
変換を必要とする．\tip{ワープ命令あたりの TLB 参照の回数は，その
アドレスが触れる異なるページの数であり，最大でワープのスレッド数である．
コアレスされたアクセスは1回か2回で済む．}したがってコアレスされない
アクセスは2度コストを払う．1度はメモリ帯域幅で（第4章），もう1度は
変換の帯域幅である．

\subsection{変換のスケーラビリティ}

GPU のアドレス変換の課題はスケーラビリティである．1つのワープ命令が
一度に多くの TLB ミスを生むことがあり，SM あたり数十のワープと数十の
SM がそれらを同時に発行する．したがってページテーブルウォーカは
高いウォークの帯域幅を維持しなければならない．\margin{Rodinia
ベンチマークとデータベースのソートで測定した値：コアレッシングと
スクラッチパッドにより変換の頻度は 6.8$\times$ 減り，TLB ミスは平均
60 個の同時ミスのバーストとして生じ，コアレッサ後の TLB は 29\,\% の
確率でミスする \cite{power2014x86}．}\cref{fig:l07-ptw-sharing}は
2つの構成を示す．SM ごとの専用 PTW では，各 SM が自分のミスをウォーク
するが，すべての SM がほとんどの時間遊んでいるウォーカの代価を払う．
共有 PTW では，SM は L1 TLB だけを持ち，そのミスを共有 L2 TLB の隣に
ある1つの高帯域幅ウォーカに送る．

\begin{figure}[tb]
  \centering
  % Private versus shared page table walkers (Lesson 7).  (a) Every SM
% has an L1 TLB and its own PTW; an L2 TLB is shared.  (b) SMs keep
% only an L1 TLB; one high-bandwidth PTW (with a page walk cache) sits
% beside the shared L2 TLB: the lecture's organisation.  The design of
% Power et al. (2014) has the shared PTW and PWC but no L2 TLB.
% Redraws the outline figures m7l1_challenges and m7l1_solutions.
% Usage: % Private versus shared page table walkers (Lesson 7).  (a) Every SM
% has an L1 TLB and its own PTW; an L2 TLB is shared.  (b) SMs keep
% only an L1 TLB; one high-bandwidth PTW (with a page walk cache) sits
% beside the shared L2 TLB: the lecture's organisation.  The design of
% Power et al. (2014) has the shared PTW and PWC but no L2 TLB.
% Redraws the outline figures m7l1_challenges and m7l1_solutions.
% Usage: % Private versus shared page table walkers (Lesson 7).  (a) Every SM
% has an L1 TLB and its own PTW; an L2 TLB is shared.  (b) SMs keep
% only an L1 TLB; one high-bandwidth PTW (with a page walk cache) sits
% beside the shared L2 TLB: the lecture's organisation.  The design of
% Power et al. (2014) has the shared PTW and PWC but no L2 TLB.
% Redraws the outline figures m7l1_challenges and m7l1_solutions.
% Usage: \input{figures/l07-ptw-sharing}
% Width: body (120 mm); the drawing is about 116 mm wide.
\begin{tikzpicture}[x=1mm, y=1mm, font=\scriptsize\sffamily,
    >={Stealth[length=1.6mm]},
    sm/.style={draw, fill=ln-box, minimum width=16mm, minimum height=12mm,
               align=center, inner sep=1pt},
    bx/.style={draw, minimum height=6mm, align=center, inner sep=1pt}]
  \foreach \ox/\title/\priv in {0/{\iflangja{(a) SM ごとの PTW}{(a) private PTW}}/1,
                                62/{\iflangja{(b) 共有 PTW}{(b) shared PTW}}/0} {
    \node[anchor=west, font=\scriptsize\sffamily\bfseries] at (\ox,38) {\title};
    \foreach \k in {0,1,2} {
      \node[sm] (sm\ox\k) at (\ox+9+\k*19,27) {};
      \node[anchor=north, inner sep=1pt] at (sm\ox\k.north) {SM\,\k};
      \ifnum\priv=1
        \node[bx, fill=ln-source!20, minimum width=7mm] at (\ox+5.2+\k*19,24.5) {\iflangja{TLB}{TLB}};
        \node[bx, fill=ln-caution!15, minimum width=7mm] at (\ox+12.8+\k*19,24.5) {PTW};
      \else
        \node[bx, fill=ln-source!20, minimum width=12mm] at (\ox+9+\k*19,24.5) {L1 TLB};
      \fi
      \draw[->] (sm\ox\k.south) -- (\ox+9+\k*19,15);
    }
    \node[bx, fill=ln-source!20, minimum width=54mm] (l2\ox) at (\ox+28,12) {\iflangja{共有 L2 TLB}{shared L2 TLB}};
    \ifnum\priv=1
      \node[bx, fill=white, minimum width=54mm] (mem\ox) at (\ox+28,0) {\iflangja{L2 キャッシュ／メモリ（ページテーブル）}{L2 cache / memory (page tables)}};
      \draw[->] (l2\ox.south) -- (l2\ox.south |- mem\ox.north);
    \else
      \node[bx, fill=white, minimum width=30mm, minimum height=8mm] (mem\ox) at (\ox+15,-1)
        {\iflangja{L2 キャッシュ／メモリ\\（ページテーブル）}{L2 cache / memory\\(page tables)}};
      \node[bx, fill=ln-caution!15, minimum width=17mm, minimum height=8mm] (pw) at (\ox+46.5,-1)
        {\iflangja{共有 PTW\\＋PWC}{shared PTW\\+ PWC}};
      \draw[->] (pw.north |- l2\ox.south) -- node[right, text=ln-muted] {\iflangja{ミス}{miss}} (pw.north);
      \draw[<->] (pw.west) -- (mem\ox.east);
    \fi
  }
\end{tikzpicture}

% Width: body (120 mm); the drawing is about 116 mm wide.
\begin{tikzpicture}[x=1mm, y=1mm, font=\scriptsize\sffamily,
    >={Stealth[length=1.6mm]},
    sm/.style={draw, fill=ln-box, minimum width=16mm, minimum height=12mm,
               align=center, inner sep=1pt},
    bx/.style={draw, minimum height=6mm, align=center, inner sep=1pt}]
  \foreach \ox/\title/\priv in {0/{\iflangja{(a) SM ごとの PTW}{(a) private PTW}}/1,
                                62/{\iflangja{(b) 共有 PTW}{(b) shared PTW}}/0} {
    \node[anchor=west, font=\scriptsize\sffamily\bfseries] at (\ox,38) {\title};
    \foreach \k in {0,1,2} {
      \node[sm] (sm\ox\k) at (\ox+9+\k*19,27) {};
      \node[anchor=north, inner sep=1pt] at (sm\ox\k.north) {SM\,\k};
      \ifnum\priv=1
        \node[bx, fill=ln-source!20, minimum width=7mm] at (\ox+5.2+\k*19,24.5) {\iflangja{TLB}{TLB}};
        \node[bx, fill=ln-caution!15, minimum width=7mm] at (\ox+12.8+\k*19,24.5) {PTW};
      \else
        \node[bx, fill=ln-source!20, minimum width=12mm] at (\ox+9+\k*19,24.5) {L1 TLB};
      \fi
      \draw[->] (sm\ox\k.south) -- (\ox+9+\k*19,15);
    }
    \node[bx, fill=ln-source!20, minimum width=54mm] (l2\ox) at (\ox+28,12) {\iflangja{共有 L2 TLB}{shared L2 TLB}};
    \ifnum\priv=1
      \node[bx, fill=white, minimum width=54mm] (mem\ox) at (\ox+28,0) {\iflangja{L2 キャッシュ／メモリ（ページテーブル）}{L2 cache / memory (page tables)}};
      \draw[->] (l2\ox.south) -- (l2\ox.south |- mem\ox.north);
    \else
      \node[bx, fill=white, minimum width=30mm, minimum height=8mm] (mem\ox) at (\ox+15,-1)
        {\iflangja{L2 キャッシュ／メモリ\\（ページテーブル）}{L2 cache / memory\\(page tables)}};
      \node[bx, fill=ln-caution!15, minimum width=17mm, minimum height=8mm] (pw) at (\ox+46.5,-1)
        {\iflangja{共有 PTW\\＋PWC}{shared PTW\\+ PWC}};
      \draw[->] (pw.north |- l2\ox.south) -- node[right, text=ln-muted] {\iflangja{ミス}{miss}} (pw.north);
      \draw[<->] (pw.west) -- (mem\ox.east);
    \fi
  }
\end{tikzpicture}

% Width: body (120 mm); the drawing is about 116 mm wide.
\begin{tikzpicture}[x=1mm, y=1mm, font=\scriptsize\sffamily,
    >={Stealth[length=1.6mm]},
    sm/.style={draw, fill=ln-box, minimum width=16mm, minimum height=12mm,
               align=center, inner sep=1pt},
    bx/.style={draw, minimum height=6mm, align=center, inner sep=1pt}]
  \foreach \ox/\title/\priv in {0/{\iflangja{(a) SM ごとの PTW}{(a) private PTW}}/1,
                                62/{\iflangja{(b) 共有 PTW}{(b) shared PTW}}/0} {
    \node[anchor=west, font=\scriptsize\sffamily\bfseries] at (\ox,38) {\title};
    \foreach \k in {0,1,2} {
      \node[sm] (sm\ox\k) at (\ox+9+\k*19,27) {};
      \node[anchor=north, inner sep=1pt] at (sm\ox\k.north) {SM\,\k};
      \ifnum\priv=1
        \node[bx, fill=ln-source!20, minimum width=7mm] at (\ox+5.2+\k*19,24.5) {\iflangja{TLB}{TLB}};
        \node[bx, fill=ln-caution!15, minimum width=7mm] at (\ox+12.8+\k*19,24.5) {PTW};
      \else
        \node[bx, fill=ln-source!20, minimum width=12mm] at (\ox+9+\k*19,24.5) {L1 TLB};
      \fi
      \draw[->] (sm\ox\k.south) -- (\ox+9+\k*19,15);
    }
    \node[bx, fill=ln-source!20, minimum width=54mm] (l2\ox) at (\ox+28,12) {\iflangja{共有 L2 TLB}{shared L2 TLB}};
    \ifnum\priv=1
      \node[bx, fill=white, minimum width=54mm] (mem\ox) at (\ox+28,0) {\iflangja{L2 キャッシュ／メモリ（ページテーブル）}{L2 cache / memory (page tables)}};
      \draw[->] (l2\ox.south) -- (l2\ox.south |- mem\ox.north);
    \else
      \node[bx, fill=white, minimum width=30mm, minimum height=8mm] (mem\ox) at (\ox+15,-1)
        {\iflangja{L2 キャッシュ／メモリ\\（ページテーブル）}{L2 cache / memory\\(page tables)}};
      \node[bx, fill=ln-caution!15, minimum width=17mm, minimum height=8mm] (pw) at (\ox+46.5,-1)
        {\iflangja{共有 PTW\\＋PWC}{shared PTW\\+ PWC}};
      \draw[->] (pw.north |- l2\ox.south) -- node[right, text=ln-muted] {\iflangja{ミス}{miss}} (pw.north);
      \draw[<->] (pw.west) -- (mem\ox.east);
    \fi
  }
\end{tikzpicture}

  \caption{(a) 各 SM の専用 PTW と共有 L2 TLB．(b) L1 TLB だけを持つ
    SM と，共有 L2 TLB の隣にあるページウォークキャッシュ（PWC）付きの
    1つの共有高帯域幅 PTW．Power らの設計は L2 TLB のない (b) である．}
  \label{fig:l07-ptw-sharing}
\end{figure}

Power，Hill と Wood は共有の構成を既知の CPU の部品から組み立てる
\cite{power2014x86}．コアレッサの後に置いた計算ユニットごとの TLB，
多数の未完了のウォークの状態を一組のバッファに保持する高度に
マルチスレッド化された PTW，そして
\term[ぺーじうぉーくきゃっしゅ]{ページウォークキャッシュ}[page walk cache]
である．ページウォークキャッシュは MMU 内の小さなキャッシュで，
ページテーブルの上位（非リーフ）段のエントリを保持し，ほとんどの
ウォークが4回未満のメモリアクセスで済むようにする．\margin{ページ
ウォークキャッシュは平均ウォークレイテンシを 95\,\% 以上削減する．
PTE がキャッシュされる GPU の L2 は 300 サイクル近くかかるので，
ウォーカにはこれが必要である．完全な設計は，\SI{4}{\kibi\byte}
ページと無修正の x86-64 ページテーブルで，平均して理想的な MMU の
2\,\% 以内に入る \cite{power2014x86}．}この設計には L2 TLB がない．
CU ごとの TLB が共有ウォーカとそのページウォークキャッシュに直接
つながる．ページウォークキャッシュの代わりに共有 L2 TLB を置くと平均で
2$\times$ 以上遅くなり，その横に追加しても性能は約 0.1\,\% しか
変わらなかった \cite[§6]{power2014x86}．

\subsection{変換のコストを減らす}

3つの技法が変換の回数またはコストを減らす．
\begin{marginfigure}
  % Segment mapping (Lesson 7): three programs, each with the virtual
% space 0x0000--0xFFFF, occupy three contiguous 64 KiB segments of
% physical memory; translation adds the segment base.
% Redraws the outline figure m7l1_solutions (left half).
% Usage: % Segment mapping (Lesson 7): three programs, each with the virtual
% space 0x0000--0xFFFF, occupy three contiguous 64 KiB segments of
% physical memory; translation adds the segment base.
% Redraws the outline figure m7l1_solutions (left half).
% Usage: % Segment mapping (Lesson 7): three programs, each with the virtual
% space 0x0000--0xFFFF, occupy three contiguous 64 KiB segments of
% physical memory; translation adds the segment base.
% Redraws the outline figure m7l1_solutions (left half).
% Usage: \input{figures/l07-segments}  (in a marginfigure)
% Width: margin (50 mm); the drawing is about 48 mm wide.
\begin{tikzpicture}[x=1mm, y=1mm, font=\scriptsize\sffamily]
  \node[anchor=south, text=ln-muted] at (8,30) {\iflangja{物理アドレス}{physical}};
  \foreach \k/\col/\lo/\hi/\prog in {0/ln-caution!20/000000/00FFFF/3,
                                     1/ln-accent!20/010000/01FFFF/2,
                                     2/ln-tip!25/020000/02FFFF/1} {
    \draw[fill=\col] (0,30-\k*10) rectangle ++(16,-10);
    \node[anchor=north west, inner sep=1pt, font=\scriptsize\ttfamily] at (0.3,30-\k*10) {\lo};
    \node[anchor=south west, inner sep=1pt, font=\scriptsize\ttfamily] at (0.3,20-\k*10) {\hi};
    \node[anchor=west, align=left] at (18,25-\k*10)
      {\iflangja{プログラム}{program}~\prog\\[-1pt]VA {\ttfamily 0000}--{\ttfamily FFFF}};
  }
\end{tikzpicture}
  (in a marginfigure)
% Width: margin (50 mm); the drawing is about 48 mm wide.
\begin{tikzpicture}[x=1mm, y=1mm, font=\scriptsize\sffamily]
  \node[anchor=south, text=ln-muted] at (8,30) {\iflangja{物理アドレス}{physical}};
  \foreach \k/\col/\lo/\hi/\prog in {0/ln-caution!20/000000/00FFFF/3,
                                     1/ln-accent!20/010000/01FFFF/2,
                                     2/ln-tip!25/020000/02FFFF/1} {
    \draw[fill=\col] (0,30-\k*10) rectangle ++(16,-10);
    \node[anchor=north west, inner sep=1pt, font=\scriptsize\ttfamily] at (0.3,30-\k*10) {\lo};
    \node[anchor=south west, inner sep=1pt, font=\scriptsize\ttfamily] at (0.3,20-\k*10) {\hi};
    \node[anchor=west, align=left] at (18,25-\k*10)
      {\iflangja{プログラム}{program}~\prog\\[-1pt]VA {\ttfamily 0000}--{\ttfamily FFFF}};
  }
\end{tikzpicture}
  (in a marginfigure)
% Width: margin (50 mm); the drawing is about 48 mm wide.
\begin{tikzpicture}[x=1mm, y=1mm, font=\scriptsize\sffamily]
  \node[anchor=south, text=ln-muted] at (8,30) {\iflangja{物理アドレス}{physical}};
  \foreach \k/\col/\lo/\hi/\prog in {0/ln-caution!20/000000/00FFFF/3,
                                     1/ln-accent!20/010000/01FFFF/2,
                                     2/ln-tip!25/020000/02FFFF/1} {
    \draw[fill=\col] (0,30-\k*10) rectangle ++(16,-10);
    \node[anchor=north west, inner sep=1pt, font=\scriptsize\ttfamily] at (0.3,30-\k*10) {\lo};
    \node[anchor=south west, inner sep=1pt, font=\scriptsize\ttfamily] at (0.3,20-\k*10) {\hi};
    \node[anchor=west, align=left] at (18,25-\k*10)
      {\iflangja{プログラム}{program}~\prog\\[-1pt]VA {\ttfamily 0000}--{\ttfamily FFFF}};
  }
\end{tikzpicture}

  \caption{セグメントマッピング：同じ仮想範囲を持つ3つのプログラムが
    3つの \SI{64}{\kibi\byte} の物理セグメントに置かれる．}
  \label{fig:l07-segments}
\end{marginfigure}
\begin{description}
  \item[セグメントマッピング．]
    \term[せぐめんとまっぴんぐ]{セグメントマッピング}[segment mapping]
    では，各プログラムの仮想空間が物理メモリの1つの連続した領域に
    対応づけられるので，変換はベースの加算になる
    （\cref{fig:l07-segments}）．図では，プログラム2の仮想アドレス
    \texttt{0x1234} は物理アドレス \texttt{0x011234} である．
    テーブルもウォークも不要だが，各プログラムの空間は物理メモリ上で
    連続していなければならない．
  \item[ラージページ．]  少数の大きなページは少ない変換で済む
    （後述）．
  \item[多段 TLB．]  SM ごとの小さな L1 TLB を，より大きな共有 L2 TLB で
    支えると（\cref{fig:l07-ptw-sharing}のように），ウォークが始まる前に，
    別の SM がすでに持ち込んだ変換で L1 のミスに応えることができる
    （ページウォークキャッシュがある場合，Power らはそこから利得を
    得られなかった．前述）．
\end{description}

\begin{definition}[ラージページ]
  \label{def:l07-large-page}
  \term[らーじぺーじ]{ラージページ}[large page]は基本ページよりはるかに
  大きなページであり，典型的には \SI{4}{\kibi\byte} の代わりに
  \SI{2}{\mebi\byte} である．オフセットは 12 ビットではなく 21 ビットで
  あるため VPN は 9 ビット短くなり，x86-64 ではウォークが1段早く終わる．
  このとき1つの TLB エントリは
  $\SI{2}{\mebi\byte}/\SI{4}{\kibi\byte} = 512$ 倍のメモリを
  カバーする．
\end{definition}

TLB が一度に対応づけられるメモリの量，すなわちエントリ数とページ
サイズの積を，その\emph{リーチ}という．\margin{CPG §19.2.2.1 Memory
Paging and Page Sizes：x86-64 の CPU は \SI{4}{\kibi\byte} の物理
ページを使うが，\enquote{NVIDIA GPUs support multiple physical page
sizes, but prefer 2MiB physical pages or larger} であり，TLB ミスは
\enquote{are generally significantly more expensive on the GPU
compared to the CPU} \cite{nvidia2024cudaguide}．}ラージページの
コストは\term[ないぶだんぺんか]{内部断片化}[internal fragmentation]で
ある．ページ単位で割り当てられたが使われないメモリであり，割り当ての
最後のページの中で失われる．

\begin{example}
  \label{ex:l07-large-pages}
  32 エントリの L1 TLB のリーチは，基本ページでは
  $32 \times \SI{4}{\kibi\byte} = \SI{128}{\kibi\byte}$，ラージページでは
  $32 \times \SI{2}{\mebi\byte} = \SI{64}{\mebi\byte}$ である．
  \SI{10}{\kibi\byte} の割り当ては \SI{4}{\kibi\byte} ページを3枚
  占めて \SI{2}{\kibi\byte} を無駄にするが，\SI{2}{\mebi\byte} ページ
  1枚では \SI{2038}{\kibi\byte}，ページの 99\,\% 以上を無駄にする
  （計算のために選んだ本書の数値）．
\end{example}

ラージページとスモールページを混在させ，メモリの各領域にそれに合った
ページサイズを与えることもできる．\sidenote{研究はこの方向で続いている．
Mosaic \cite{ausavarungnirun2017mosaic} はアプリケーションを変更せずに
複数のページサイズを同時にサポートし，複数アプリケーションの
ワークロードで平均 55.5\,\%（同種）と 29.7\,\%（異種）の向上を報告して
おり，理想的な TLB との差はそれぞれ 6.8\,\% と 15.4\,\% 以内である．
複数の GPU については，Li ら \cite{li2021tlb} が変換をピアの GPU 間で
共有し，TLB エントリを IOMMU の TLB にあふれさせることで，単一および
複数アプリケーションのワークロードで平均 23.5\,\% と 16.3\,\% の向上を
得ている．}

\begin{keypoint}
  GPU はワープ全体のアドレスを変換する．命令が必要とする変換の数は
  コアレッシングが決め，コアレスされないアクセスはスレッドごとに1回の
  変換を必要としうるし，ウォーカは多数のミスを同時に処理しなければ
  ならない．ページウォークキャッシュを持つ共有の高帯域幅ウォーカ，
  多段 TLB，セグメント，ラージページがコストを減らす．ラージページは
  TLB のリーチと内部断片化を交換する．
\end{keypoint}

\section{統合仮想アドレス空間とオンデマンドページング}
\label{sec:l07-uva}

仮想メモリの第2の側面はページの割り当てであり，GPU が CPU と最も
異なるのはここである．\source{M7，UVA とオンデマンドページング；
CPG §19 Unified Memory Programming；BPG §10.1.4 Unified Virtual
Addressing；\cite{ganguly2019interplay,kim2020batch}．}本節では明示的な
コピーと共有アドレス空間を比較し，GPU のページフォールトをホストまで
たどって戻り，オンデマンドページングを実用的なコストにする
プリフェッチと追い出しのポリシを紹介する．

\subsection{コピー後実行と共有アドレス空間}

\begin{definition}[コピー後実行モデル]
  \label{def:l07-copy-then-execute}
  \term[こぴーごじっこうもでる]{コピー後実行モデル}[copy-then-execute model]
  では，ホストはカーネルを起動する前にデバイスメモリを割り当て
  （\texttt{cudaMalloc}），入力をそこへコピーし（\texttt{cudaMemcpy}），
  終了後に結果をコピーして戻す．カーネルが触れるすべてのデータは，
  カーネルの開始前にデバイスメモリにある．
\end{definition}

データ転送は明示的であり，カーネル起動前の準備の中で行われる．
これには1つの長所がある．プログラマがスケジュールできる一括転送は，
第5章のピン留めバッファとストリームが示したように，PCIe の帯域幅を
使い切る．短所は，プログラマがすべてのバッファとすべてのコピーを
手作業で管理することであり，ポインタを含むデータ構造はまったく共有
できないことである．

\caution{\enquote{SVM} は OpenCL 2.0 の用語である．CUDA はこれを
使わず，共有アドレス空間を統合仮想アドレッシング，移動する変種を
統合（マネージド）メモリと呼ぶ．}%
\begin{definition}[共有仮想メモリ]
  \label{def:l07-svm}
  OpenCL/CUDA のプログラミングモデルで用いられる名前である
  \term[きょうゆうかそうめもり]{共有仮想メモリ}[shared virtual memory (SVM)]（SVM）
  では，CPU と GPU は1つの仮想アドレス空間を使うので，1つのポインタは
  両方で同じデータを意味する（\cref{fig:l07-uva}）．\index{SVM|see{共有仮想メモリ}}
\end{definition}

\begin{figure}[tb]
  \centering
  % Separate memory spaces versus a unified virtual address space
% (Lesson 7).  (a) The CPU writes and reads a buffer in CPU memory, the
% GPU one in GPU memory, and explicit copies move data between them.
% (b) One logical buffer in a unified space is read and written by
% both processors; no copy appears in the program.
% Redraws the outline figure m7l2_uva.
% Usage: % Separate memory spaces versus a unified virtual address space
% (Lesson 7).  (a) The CPU writes and reads a buffer in CPU memory, the
% GPU one in GPU memory, and explicit copies move data between them.
% (b) One logical buffer in a unified space is read and written by
% both processors; no copy appears in the program.
% Redraws the outline figure m7l2_uva.
% Usage: % Separate memory spaces versus a unified virtual address space
% (Lesson 7).  (a) The CPU writes and reads a buffer in CPU memory, the
% GPU one in GPU memory, and explicit copies move data between them.
% (b) One logical buffer in a unified space is read and written by
% both processors; no copy appears in the program.
% Redraws the outline figure m7l2_uva.
% Usage: \input{figures/l07-uva}
% Width: body (120 mm); the drawing is about 112 mm wide.
\begin{tikzpicture}[x=1mm, y=1mm, font=\scriptsize\sffamily,
    >={Stealth[length=1.6mm]},
    pr/.style={draw, fill=ln-accent!20, minimum width=12mm, minimum height=6mm},
    buf/.style={draw, fill=ln-tip!25, minimum width=14mm, minimum height=6mm}]
  % (a) separate
  \node[anchor=west, font=\scriptsize\sffamily\bfseries] at (0,33)
    {\iflangja{(a) 分離したメモリ空間}{(a) separate memory spaces}};
  \node[pr] (c1) at (9,24) {CPU};
  \node[pr] (g1) at (45,24) {GPU};
  \draw[fill=ln-box] (0,15) rectangle (20,3);
  \node[anchor=north, text=ln-muted] at (10,15) {\iflangja{CPU メモリ}{CPU memory}};
  \node[buf] (b1) at (10,7) {\iflangja{バッファ}{buffer}};
  \draw[fill=ln-box] (34,15) rectangle (54,3);
  \node[anchor=north, text=ln-muted] at (44,15) {\iflangja{GPU メモリ}{GPU memory}};
  \node[buf] (b2) at (44,7) {\iflangja{バッファ}{buffer}};
  \draw[<->] (c1) -- (c1 |- 0,15);
  \draw[<->] (g1) -- (g1 |- 0,15);
  \draw[->, thick, ln-caution] (b1.east) -- node[above, text=ln-caution] {\iflangja{コピー}{copy}} (b2.west);
  % (b) unified
  \node[anchor=west, font=\scriptsize\sffamily\bfseries] at (62,33)
    {\iflangja{(b) 統合仮想アドレス空間}{(b) unified virtual address space}};
  \node[pr] (c2) at (71,24) {CPU};
  \node[pr] (g2) at (103,24) {GPU};
  \draw[fill=ln-box] (62,15) rectangle (112,3);
  \node[anchor=north, text=ln-muted] at (87,2.5) {\iflangja{統合メモリ空間}{unified memory space}};
  \node[buf] (b3) at (87,7) {\iflangja{バッファ}{buffer}};
  \draw[<->] (c2) -- (b3);
  \draw[<->] (g2) -- (b3);
\end{tikzpicture}

% Width: body (120 mm); the drawing is about 112 mm wide.
\begin{tikzpicture}[x=1mm, y=1mm, font=\scriptsize\sffamily,
    >={Stealth[length=1.6mm]},
    pr/.style={draw, fill=ln-accent!20, minimum width=12mm, minimum height=6mm},
    buf/.style={draw, fill=ln-tip!25, minimum width=14mm, minimum height=6mm}]
  % (a) separate
  \node[anchor=west, font=\scriptsize\sffamily\bfseries] at (0,33)
    {\iflangja{(a) 分離したメモリ空間}{(a) separate memory spaces}};
  \node[pr] (c1) at (9,24) {CPU};
  \node[pr] (g1) at (45,24) {GPU};
  \draw[fill=ln-box] (0,15) rectangle (20,3);
  \node[anchor=north, text=ln-muted] at (10,15) {\iflangja{CPU メモリ}{CPU memory}};
  \node[buf] (b1) at (10,7) {\iflangja{バッファ}{buffer}};
  \draw[fill=ln-box] (34,15) rectangle (54,3);
  \node[anchor=north, text=ln-muted] at (44,15) {\iflangja{GPU メモリ}{GPU memory}};
  \node[buf] (b2) at (44,7) {\iflangja{バッファ}{buffer}};
  \draw[<->] (c1) -- (c1 |- 0,15);
  \draw[<->] (g1) -- (g1 |- 0,15);
  \draw[->, thick, ln-caution] (b1.east) -- node[above, text=ln-caution] {\iflangja{コピー}{copy}} (b2.west);
  % (b) unified
  \node[anchor=west, font=\scriptsize\sffamily\bfseries] at (62,33)
    {\iflangja{(b) 統合仮想アドレス空間}{(b) unified virtual address space}};
  \node[pr] (c2) at (71,24) {CPU};
  \node[pr] (g2) at (103,24) {GPU};
  \draw[fill=ln-box] (62,15) rectangle (112,3);
  \node[anchor=north, text=ln-muted] at (87,2.5) {\iflangja{統合メモリ空間}{unified memory space}};
  \node[buf] (b3) at (87,7) {\iflangja{バッファ}{buffer}};
  \draw[<->] (c2) -- (b3);
  \draw[<->] (g2) -- (b3);
\end{tikzpicture}

% Width: body (120 mm); the drawing is about 112 mm wide.
\begin{tikzpicture}[x=1mm, y=1mm, font=\scriptsize\sffamily,
    >={Stealth[length=1.6mm]},
    pr/.style={draw, fill=ln-accent!20, minimum width=12mm, minimum height=6mm},
    buf/.style={draw, fill=ln-tip!25, minimum width=14mm, minimum height=6mm}]
  % (a) separate
  \node[anchor=west, font=\scriptsize\sffamily\bfseries] at (0,33)
    {\iflangja{(a) 分離したメモリ空間}{(a) separate memory spaces}};
  \node[pr] (c1) at (9,24) {CPU};
  \node[pr] (g1) at (45,24) {GPU};
  \draw[fill=ln-box] (0,15) rectangle (20,3);
  \node[anchor=north, text=ln-muted] at (10,15) {\iflangja{CPU メモリ}{CPU memory}};
  \node[buf] (b1) at (10,7) {\iflangja{バッファ}{buffer}};
  \draw[fill=ln-box] (34,15) rectangle (54,3);
  \node[anchor=north, text=ln-muted] at (44,15) {\iflangja{GPU メモリ}{GPU memory}};
  \node[buf] (b2) at (44,7) {\iflangja{バッファ}{buffer}};
  \draw[<->] (c1) -- (c1 |- 0,15);
  \draw[<->] (g1) -- (g1 |- 0,15);
  \draw[->, thick, ln-caution] (b1.east) -- node[above, text=ln-caution] {\iflangja{コピー}{copy}} (b2.west);
  % (b) unified
  \node[anchor=west, font=\scriptsize\sffamily\bfseries] at (62,33)
    {\iflangja{(b) 統合仮想アドレス空間}{(b) unified virtual address space}};
  \node[pr] (c2) at (71,24) {CPU};
  \node[pr] (g2) at (103,24) {GPU};
  \draw[fill=ln-box] (62,15) rectangle (112,3);
  \node[anchor=north, text=ln-muted] at (87,2.5) {\iflangja{統合メモリ空間}{unified memory space}};
  \node[buf] (b3) at (87,7) {\iflangja{バッファ}{buffer}};
  \draw[<->] (c2) -- (b3);
  \draw[<->] (g2) -- (b3);
\end{tikzpicture}

  \caption{(a) 分離したメモリ空間：各プロセッサは自分のバッファを
    読み書きし，明示的なコピーがデータを動かす．(b) 統合仮想アドレス
    空間：両方のプロセッサが読み書きする1つの論理バッファ．}
  \label{fig:l07-uva}
\end{figure}

第5章の\term*{統合仮想アドレッシング（UVA）}は分離したメモリの上の
1つのアドレス空間である．ページも自動的に移動するようになると，共有
アドレス空間はプログラマビリティをさらに高める．\margin{第3章は
\texttt{cudaMallocManaged} のプログラミング面を紹介した．第5章は
UVA（1つのアドレス空間）とユニファイドメモリ（自動的なページ移動）を区別した．}
高性能計算の GPU のほとんどはディスクリートである．すなわち自分の
物理メモリを持つので，共有アドレス空間でもホストメモリとデバイス
メモリの間のコピーは依然として必要であり，それが今度は要求に応じて
行われる．これには\emph{動的なページ割り当て}が必要である．ページを
どこに置くかの決定が，カーネルの実行中に下される．

\begin{definition}[オンデマンドページング]
  \label{def:l07-on-demand}
  \term[おんでまんどぺーじんぐ]{オンデマンドページング}[on-demand paging]
  のもとでは，マネージドメモリ（\texttt{cudaMallocManaged} で割り当てた
  メモリ）は前もってデバイスメモリに置かれない．常駐していないページへの GPU の最初の
  アクセスがフォールトを起こし，そのときページがデバイスメモリに
  割り当てられ，そのデータがホストから移される．
\end{definition}

$n$ 個の float をその場で何倍かするカーネル \texttt{scale} について，
2つのモデルを並べる．

\begin{lstlisting}[language=C++, numbers=none]
// コピー後実行
float *d;
cudaMalloc(&d, n * sizeof(float));
cudaMemcpy(d, h, n * sizeof(float),
           cudaMemcpyHostToDevice);
scale<<<(n + 255) / 256, 256>>>(d, n);
cudaMemcpy(h, d, n * sizeof(float),
           cudaMemcpyDeviceToHost);

// オンデマンドページング
float *m;
cudaMallocManaged(&m, n * sizeof(float));
for (int i = 0; i < n; ++i) m[i] = i;  // CPU が触れる
scale<<<(n + 255) / 256, 256>>>(m, n);  // GPU がフォールト
cudaDeviceSynchronize();      // m は CPU で有効
\end{lstlisting}

\margin{CPG §19.3.2.1 Data Migration and Coherency：Compute
Capability 6.x より前は，カーネル起動のたびにすべてのマネージドメモリが
GPU へ移された．Compute Capability 6.x（Pascal）は GPU のページ
フォールトを導入し，ページを要求に応じて移すようになった
\cite{nvidia2024cudaguide}．}%
2つ目の版にはプログラム中のコピーがないが，コピーはやはり起こる．各ページは
カーネルが初めて触れたときに PCIe を渡る．

\subsection{ページフォールトの経路}

\begin{definition}[ファーフォールトと GMMU]
  \label{def:l07-far-fault}
  ユニファイドメモリを持つ GPU では，
  \term[じーぴーゆーめもりかんりゆにっと]{GPUメモリ管理ユニット}[GPU memory management unit (GMMU)]（GMMU）
  が SM の TLB ミスを受け取り，GPU のページテーブルをウォークする．
  \term[ふぁーふぉーると]{ファーフォールト}[far fault]とは，デバイス
  メモリに存在しないページに対する GPU のページフォールトであり，
  ホストのドライバがページを割り当ててデータを移すことで解決しなければ
  ならない．これに対し\emph{ニアフォールト}は，デバイスメモリが応える
  L2 ミスである．\index{GMMU|see{GPUメモリ管理ユニット}}
\end{definition}

\cref{fig:l07-fault-path}はファーフォールトをたどる．
\begin{enumerate}
  \item GMMU のウォークは有効な PTE を見つけられない．フォールトは
    GPU のページフォールト用 MSHR（miss status holding register）に
    記録され，フォールトしたワープは待つ．
  \item フォールトは PCIe を越えて CPU へ送られる．
  \item オペレーティングシステムと GPU ドライバがそれを処理する．
    デバイスメモリにページを割り当て，データを移す．
  \item 新しい対応が送り返され，PTE が有効にされ，アクセスが再実行
    される．
\end{enumerate}
この往復には数十マイクロ秒かかる．\caution{講義 vs 論文：出典の
一節はこのレイテンシを\enquote{tens of milliseconds}としているが，
正しくは数十マイクロ秒である．先行研究はフォールトあたり
\SIrange{20}{50}{\micro\second} を報告しており，Kim らはフォールトの
バッチ全体について \SIrange{223}{553}{\micro\second}（中央値
\SI{313}{\micro\second}）を測定している \cite{kim2020batch}．}

\begin{figure}[b]
  \centering
  % Far-fault handling with on-demand paging on a discrete GPU
% (Lesson 7).  (1) A TLB miss goes to the GMMU, whose walk finds no
% valid PTE: a far fault, recorded in the far-fault MSHRs.  (2) The
% fault is sent over PCIe to the host.  (3) The GPU driver and the OS
% allocate a page and migrate the data.  (4) The new mapping returns,
% the PTE is validated and the access is replayed.
% Redraws the outline figures m7l2_page_fault_handling and
% m7l2_challenge_ondemand.
% Usage: % Far-fault handling with on-demand paging on a discrete GPU
% (Lesson 7).  (1) A TLB miss goes to the GMMU, whose walk finds no
% valid PTE: a far fault, recorded in the far-fault MSHRs.  (2) The
% fault is sent over PCIe to the host.  (3) The GPU driver and the OS
% allocate a page and migrate the data.  (4) The new mapping returns,
% the PTE is validated and the access is replayed.
% Redraws the outline figures m7l2_page_fault_handling and
% m7l2_challenge_ondemand.
% Usage: % Far-fault handling with on-demand paging on a discrete GPU
% (Lesson 7).  (1) A TLB miss goes to the GMMU, whose walk finds no
% valid PTE: a far fault, recorded in the far-fault MSHRs.  (2) The
% fault is sent over PCIe to the host.  (3) The GPU driver and the OS
% allocate a page and migrate the data.  (4) The new mapping returns,
% the PTE is validated and the access is replayed.
% Redraws the outline figures m7l2_page_fault_handling and
% m7l2_challenge_ondemand.
% Usage: \input{figures/l07-fault-path}
% Width: body (120 mm); the drawing is about 116 mm wide.
\begin{tikzpicture}[x=1mm, y=1mm, font=\scriptsize\sffamily,
    >={Stealth[length=1.6mm]},
    bx/.style={draw, minimum height=6mm, align=center, inner sep=1pt},
    stp/.style={circle, draw, fill=white, inner sep=0.6pt, font=\scriptsize\sffamily\bfseries}]
  % GPU
  \draw[fill=ln-box] (0,0) rectangle (56,40);
  \node[anchor=north west, font=\scriptsize\sffamily\bfseries] at (0.5,39.5) {GPU};
  \foreach \k in {0,1,2} {
    \node[bx, fill=ln-accent!20, minimum width=14mm] at (8+\k*17,31) {SM\,\k\ +\,TLB};
  }
  \node[bx, fill=ln-caution!15, minimum width=50mm, minimum height=13mm] (gmmu) at (28,17)
    {};
  \node[anchor=north, inner sep=1pt] at (gmmu.north) {GMMU};
  \node[bx, fill=white, minimum width=14mm] at (11,15) {L2 TLB};
  \node[bx, fill=white, minimum width=14mm] at (28,15) {\iflangja{ウォーク}{walk}};
  \node[bx, fill=white, minimum width=14mm] at (45,15) {\iflangja{フォールト\\MSHR}{fault\\MSHRs}};
  \node[bx, fill=ln-tip!25, minimum width=50mm] (dmem) at (28,4.5) {\iflangja{デバイスメモリ（ページテーブル）}{device memory (page table)}};
  \foreach \k in {0,1,2} { \draw[->] (8+\k*17,28) -- (8+\k*17,23.5); }
  % PCIe
  \draw[fill=ln-muted!15] (60,14) rectangle (72,26);
  \node[rotate=90] at (66,20) {PCIe};
  % CPU
  \draw[fill=ln-box] (76,0) rectangle (116,40);
  \node[anchor=north west, font=\scriptsize\sffamily\bfseries] at (76.5,39.5) {CPU};
  \node[bx, fill=ln-accent!20, minimum width=34mm] (drv) at (96,26) {\iflangja{OS ＋ GPU ドライバ／\\ランタイム}{OS + GPU driver\\and runtime}};
  \node[bx, fill=ln-tip!25, minimum width=34mm] (hmem) at (96,7) {\iflangja{ホストメモリ}{host memory}};
  \draw[<->] (drv) -- (hmem);
  % fault and reply
  \draw[->, thick, ln-caution] (53,24) -- (80,29);
  \node[stp] at (49,25) {1};
  \node[stp] at (66,30) {2};
  \node[anchor=south, text=ln-caution, align=center] at (66,31.8) {\iflangja{ページフォールト}{page fault}};
  \node[stp] at (106,17) {3};
  \draw[->, thick, ln-tip] (80,22) -- (53,11);
  \node[stp] at (66,10) {4};
  \node[anchor=north, text=ln-tip, align=center] at (66,8.5) {\iflangja{新しいマッピング\\とデータ}{new mapping\\and data}};
\end{tikzpicture}

% Width: body (120 mm); the drawing is about 116 mm wide.
\begin{tikzpicture}[x=1mm, y=1mm, font=\scriptsize\sffamily,
    >={Stealth[length=1.6mm]},
    bx/.style={draw, minimum height=6mm, align=center, inner sep=1pt},
    stp/.style={circle, draw, fill=white, inner sep=0.6pt, font=\scriptsize\sffamily\bfseries}]
  % GPU
  \draw[fill=ln-box] (0,0) rectangle (56,40);
  \node[anchor=north west, font=\scriptsize\sffamily\bfseries] at (0.5,39.5) {GPU};
  \foreach \k in {0,1,2} {
    \node[bx, fill=ln-accent!20, minimum width=14mm] at (8+\k*17,31) {SM\,\k\ +\,TLB};
  }
  \node[bx, fill=ln-caution!15, minimum width=50mm, minimum height=13mm] (gmmu) at (28,17)
    {};
  \node[anchor=north, inner sep=1pt] at (gmmu.north) {GMMU};
  \node[bx, fill=white, minimum width=14mm] at (11,15) {L2 TLB};
  \node[bx, fill=white, minimum width=14mm] at (28,15) {\iflangja{ウォーク}{walk}};
  \node[bx, fill=white, minimum width=14mm] at (45,15) {\iflangja{フォールト\\MSHR}{fault\\MSHRs}};
  \node[bx, fill=ln-tip!25, minimum width=50mm] (dmem) at (28,4.5) {\iflangja{デバイスメモリ（ページテーブル）}{device memory (page table)}};
  \foreach \k in {0,1,2} { \draw[->] (8+\k*17,28) -- (8+\k*17,23.5); }
  % PCIe
  \draw[fill=ln-muted!15] (60,14) rectangle (72,26);
  \node[rotate=90] at (66,20) {PCIe};
  % CPU
  \draw[fill=ln-box] (76,0) rectangle (116,40);
  \node[anchor=north west, font=\scriptsize\sffamily\bfseries] at (76.5,39.5) {CPU};
  \node[bx, fill=ln-accent!20, minimum width=34mm] (drv) at (96,26) {\iflangja{OS ＋ GPU ドライバ／\\ランタイム}{OS + GPU driver\\and runtime}};
  \node[bx, fill=ln-tip!25, minimum width=34mm] (hmem) at (96,7) {\iflangja{ホストメモリ}{host memory}};
  \draw[<->] (drv) -- (hmem);
  % fault and reply
  \draw[->, thick, ln-caution] (53,24) -- (80,29);
  \node[stp] at (49,25) {1};
  \node[stp] at (66,30) {2};
  \node[anchor=south, text=ln-caution, align=center] at (66,31.8) {\iflangja{ページフォールト}{page fault}};
  \node[stp] at (106,17) {3};
  \draw[->, thick, ln-tip] (80,22) -- (53,11);
  \node[stp] at (66,10) {4};
  \node[anchor=north, text=ln-tip, align=center] at (66,8.5) {\iflangja{新しいマッピング\\とデータ}{new mapping\\and data}};
\end{tikzpicture}

% Width: body (120 mm); the drawing is about 116 mm wide.
\begin{tikzpicture}[x=1mm, y=1mm, font=\scriptsize\sffamily,
    >={Stealth[length=1.6mm]},
    bx/.style={draw, minimum height=6mm, align=center, inner sep=1pt},
    stp/.style={circle, draw, fill=white, inner sep=0.6pt, font=\scriptsize\sffamily\bfseries}]
  % GPU
  \draw[fill=ln-box] (0,0) rectangle (56,40);
  \node[anchor=north west, font=\scriptsize\sffamily\bfseries] at (0.5,39.5) {GPU};
  \foreach \k in {0,1,2} {
    \node[bx, fill=ln-accent!20, minimum width=14mm] at (8+\k*17,31) {SM\,\k\ +\,TLB};
  }
  \node[bx, fill=ln-caution!15, minimum width=50mm, minimum height=13mm] (gmmu) at (28,17)
    {};
  \node[anchor=north, inner sep=1pt] at (gmmu.north) {GMMU};
  \node[bx, fill=white, minimum width=14mm] at (11,15) {L2 TLB};
  \node[bx, fill=white, minimum width=14mm] at (28,15) {\iflangja{ウォーク}{walk}};
  \node[bx, fill=white, minimum width=14mm] at (45,15) {\iflangja{フォールト\\MSHR}{fault\\MSHRs}};
  \node[bx, fill=ln-tip!25, minimum width=50mm] (dmem) at (28,4.5) {\iflangja{デバイスメモリ（ページテーブル）}{device memory (page table)}};
  \foreach \k in {0,1,2} { \draw[->] (8+\k*17,28) -- (8+\k*17,23.5); }
  % PCIe
  \draw[fill=ln-muted!15] (60,14) rectangle (72,26);
  \node[rotate=90] at (66,20) {PCIe};
  % CPU
  \draw[fill=ln-box] (76,0) rectangle (116,40);
  \node[anchor=north west, font=\scriptsize\sffamily\bfseries] at (76.5,39.5) {CPU};
  \node[bx, fill=ln-accent!20, minimum width=34mm] (drv) at (96,26) {\iflangja{OS ＋ GPU ドライバ／\\ランタイム}{OS + GPU driver\\and runtime}};
  \node[bx, fill=ln-tip!25, minimum width=34mm] (hmem) at (96,7) {\iflangja{ホストメモリ}{host memory}};
  \draw[<->] (drv) -- (hmem);
  % fault and reply
  \draw[->, thick, ln-caution] (53,24) -- (80,29);
  \node[stp] at (49,25) {1};
  \node[stp] at (66,30) {2};
  \node[anchor=south, text=ln-caution, align=center] at (66,31.8) {\iflangja{ページフォールト}{page fault}};
  \node[stp] at (106,17) {3};
  \draw[->, thick, ln-tip] (80,22) -- (53,11);
  \node[stp] at (66,10) {4};
  \node[anchor=north, text=ln-tip, align=center] at (66,8.5) {\iflangja{新しいマッピング\\とデータ}{new mapping\\and data}};
\end{tikzpicture}

  \caption{ディスクリート GPU でのファーフォールト：(1) GMMU が有効な
    PTE を見つけられずフォールトを記録する．(2) フォールトが PCIe を
    渡る．(3) OS と GPU ドライバがページを割り当てて移す．(4) 対応が
    戻り，アクセスが再実行される．}
  \label{fig:l07-fault-path}
\end{figure}

両方のプロセッサは，デバイス用の変換ユニットを通じてメモリに到達する．

\begin{definition}[IOMMU と DMA]
  \label{def:l07-iommu}
  \term[でぃーえむえー]{DMA}[DMA (direct memory access)]（direct
  memory access）とは，CPU がコピーを実行することなく，デバイスと
  メモリの間でデータを転送することである．
  \term[あいおーえむえむゆー]{IOMMU}[IOMMU]（input--output memory
  management unit）は，デバイスが DMA に用いるアドレスを変換し，
  デバイスが物理アドレスではなく仮想アドレスで動作できるようにする
  （I/O 仮想化）．\index{direct memory access|see{DMA}}
\end{definition}

\margin{Ganguly らは，Pascal GPU でのファーフォールトの処理
レイテンシを，データの PCIe 転送の前で典型的に
\SI{45}{\micro\second} としている \cite{ganguly2019interplay}．}%
GPU のメモリはかつて，周辺デバイスのメモリとして IOMMU を通じてのみ
管理されていた．共有アドレス空間では CPU があらゆる対応の変更に関与し，
これがオンデマンドの GPU メモリのコストである．ページ管理はホストが
行い，CPU から対応を取ってくるには PCIe の往復に加えて CPU の割り込み
処理ルーチンがかかり，しばしば数十マイクロ秒になる．\sidenote{ドライバは
フォールトを1つずつ処理するのではない．フォールトをバッファに集めて
バッチで処理し，バッチの処理中に起きたフォールトは次のバッチを待つ．
Kim ら \cite{kim2020batch} はこの直列化がスループットを損なうことを示し，
より大きなバッチ（スレッドブロックのコンテキストを切り替えて，より多くの
フォールトがまとめて起こるようにする）と，移動と重ね合わせた追い出しを
提案している．最先端のプリフェッチに対して平均 2$\times$ の高速化，
バッチは 2.27$\times$ 大きく 51\,\% 少なく，バッチ処理時間は 27\,\%
短い．}3つの対策が考えられる．
\begin{itemize}
  \item \emph{ページをプリフェッチする}：必要になるページを予測し，
    フォールトする前に移す．
  \item \emph{ページ対応の管理を GPU 側に移す}：フォールトがホストを
    必要としないようにする．
  \item \emph{CPU と GPU を協調させる}：CPU はメモリの大きなチャンクを
    管理し，GPU は与えられたチャンク内の小さなページを管理する．
\end{itemize}

NVIDIA のドライバが実装している対策はプリフェッチである
\cite{ganguly2019interplay}．

\subsection{ページのプリフェッチ}

\begin{definition}[プリフェッチ]
  \label{def:l07-prefetching}
  将来のアドレスの予測に基づいて，要求される前にデータを取ってくる
  ことを\term[ぷりふぇっち]{プリフェッチ}[prefetching]という．
  これを統合アドレス空間のページに適用したものが
  \term[ぺーじぷりふぇっち]{ページプリフェッチ}[page prefetching]で
  あり，GPU がフォールトする前にページをデバイスメモリへ移す．
\end{definition}

アドレス \texttt{0x40}，\texttt{0x80}，\texttt{0xC0} を見た逐次予測器は，
次に \texttt{0x100} と \texttt{0x140} を予測する．連続するインデックスで
\texttt{a[idx]} と \texttt{b[idx]} を読むベクトル加算はまさにそのような
系列を生み，そのような配列のページアドレスも同じく予測可能である．
逐次ページ予測器はページ $p$ がフォールトしたときにページ $p+1$ を
取ってくる．

\begin{marginfigure}
  % Page prefetching under out-of-order thread blocks (Lesson 7): an
% array spans nine pages; block 0 reads pages 0--2, block 1 pages
% 3--5, block 2 pages 6--8.  The blocks run in the order 2, 0, 1, so
% the faults arrive as 6 7 8 0 1 2 3 4 5 and a predictor of "next
% page" mispredicts at every block boundary.
% Redraws the outline figure m7l2_preteching (page numbers made
% consistent).
% Usage: % Page prefetching under out-of-order thread blocks (Lesson 7): an
% array spans nine pages; block 0 reads pages 0--2, block 1 pages
% 3--5, block 2 pages 6--8.  The blocks run in the order 2, 0, 1, so
% the faults arrive as 6 7 8 0 1 2 3 4 5 and a predictor of "next
% page" mispredicts at every block boundary.
% Redraws the outline figure m7l2_preteching (page numbers made
% consistent).
% Usage: % Page prefetching under out-of-order thread blocks (Lesson 7): an
% array spans nine pages; block 0 reads pages 0--2, block 1 pages
% 3--5, block 2 pages 6--8.  The blocks run in the order 2, 0, 1, so
% the faults arrive as 6 7 8 0 1 2 3 4 5 and a predictor of "next
% page" mispredicts at every block boundary.
% Redraws the outline figure m7l2_preteching (page numbers made
% consistent).
% Usage: \input{figures/l07-block-pages}  (in a marginfigure)
% Width: margin (50 mm); the drawing is about 48 mm wide.
\begin{tikzpicture}[x=1mm, y=1mm, font=\scriptsize\sffamily]
  \foreach \p in {0,...,8} {
    \pgfmathtruncatemacro{\b}{\p/3}
    \ifcase\b \def\col{ln-accent!25}\or \def\col{ln-caution!20}\or \def\col{ln-tip!25}\fi
    \draw[fill=\col] (\p*5.2,10) rectangle ++(5.2,6) node[midway] {\p};
  }
  \node[anchor=south west, text=ln-muted] at (0,16) {\iflangja{ページ}{page}};
  \foreach \b/\col in {0/ln-accent!25, 1/ln-caution!20, 2/ln-tip!25} {
    \draw[fill=\col] (\b*15.6+0.5,2) rectangle ++(14.6,5)
      node[midway] {\iflangja{ブロック}{block}~\b};
  }
  \node[anchor=north west, text=ln-muted, align=left] at (0,1)
    {\iflangja{実行順：2, 0, 1}{run order: 2, 0, 1}};
\end{tikzpicture}
  (in a marginfigure)
% Width: margin (50 mm); the drawing is about 48 mm wide.
\begin{tikzpicture}[x=1mm, y=1mm, font=\scriptsize\sffamily]
  \foreach \p in {0,...,8} {
    \pgfmathtruncatemacro{\b}{\p/3}
    \ifcase\b \def\col{ln-accent!25}\or \def\col{ln-caution!20}\or \def\col{ln-tip!25}\fi
    \draw[fill=\col] (\p*5.2,10) rectangle ++(5.2,6) node[midway] {\p};
  }
  \node[anchor=south west, text=ln-muted] at (0,16) {\iflangja{ページ}{page}};
  \foreach \b/\col in {0/ln-accent!25, 1/ln-caution!20, 2/ln-tip!25} {
    \draw[fill=\col] (\b*15.6+0.5,2) rectangle ++(14.6,5)
      node[midway] {\iflangja{ブロック}{block}~\b};
  }
  \node[anchor=north west, text=ln-muted, align=left] at (0,1)
    {\iflangja{実行順：2, 0, 1}{run order: 2, 0, 1}};
\end{tikzpicture}
  (in a marginfigure)
% Width: margin (50 mm); the drawing is about 48 mm wide.
\begin{tikzpicture}[x=1mm, y=1mm, font=\scriptsize\sffamily]
  \foreach \p in {0,...,8} {
    \pgfmathtruncatemacro{\b}{\p/3}
    \ifcase\b \def\col{ln-accent!25}\or \def\col{ln-caution!20}\or \def\col{ln-tip!25}\fi
    \draw[fill=\col] (\p*5.2,10) rectangle ++(5.2,6) node[midway] {\p};
  }
  \node[anchor=south west, text=ln-muted] at (0,16) {\iflangja{ページ}{page}};
  \foreach \b/\col in {0/ln-accent!25, 1/ln-caution!20, 2/ln-tip!25} {
    \draw[fill=\col] (\b*15.6+0.5,2) rectangle ++(14.6,5)
      node[midway] {\iflangja{ブロック}{block}~\b};
  }
  \node[anchor=north west, text=ln-muted, align=left] at (0,1)
    {\iflangja{実行順：2, 0, 1}{run order: 2, 0, 1}};
\end{tikzpicture}

  \caption{2, 0, 1 の順に実行される3つのスレッドブロックが読む9ページ：
    フォールトは 6, 7, 8, 0, 1, \ldots の順に届く．}
  \label{fig:l07-block-pages}
\end{marginfigure}

GPU ではこの予測は破綻する．スレッドブロックは順序どおりに実行され
ないからである（第3章）．\cref{fig:l07-block-pages}では各ブロックが
連続する3ページを読むが，ブロックは 2, 0, 1 の順に実行され，
フォールトは 6, 7, 8, 0, 1, 2, 3, 4, 5 の順に届く．ページ8の後，
\enquote{次のページ}予測器はどのブロックも読まないページ9を取ってきて
しまい，ページ0への後戻りは予見できないフォールトになる．空間的局所性は
フォールトの近傍には存在するが，フォールトの順序には存在しない．

\begin{definition}[ツリーベース近傍プリフェッチ]
  \label{def:l07-tbnp}
  NVIDIA のユニファイドメモリドライバのポリシは
  \term[つりーべーすきんぼうぷりふぇっち]{ツリーベース近傍プリフェッチ}[tree-based neighborhood prefetching]
  （TBNp）と呼ばれ，マネージドな割り当ての各 \SI{2}{\mebi\byte} 領域を
  2つの \SI{1}{\mebi\byte} の半分に，それらをさらに
  \SI{64}{\kibi\byte} のノードに分割する．\SI{4}{\kibi\byte} の
  フォールトは \SI{64}{\kibi\byte} のノード全体，16ページを移す．
  1つの親の下のいくつかのノードがフォールトすると，最大
  \SI{1}{\mebi\byte} の親ノードが取ってこられる．\index{TBNp|see{ツリーベース近傍プリフェッチ}}
\end{definition}

このツリーは，フォールトの順序にかかわらず各フォールトの周囲の空間的
局所性をとらえ，1回のプリフェッチの大きさを抑える．\caution{講義 vs
論文：Ganguly ら \cite{ganguly2019interplay} は正確な規則を
GTX 1080 Ti で測定した．ツリーは \SI{64}{\kibi\byte} の葉を持つ完全
二分木であり，その間に 128，256，\SI{512}{\kibi\byte} の段がある．
ノードは\enquote{to-be valid size}が容量の\enquote{strictly greater
than 50\,\%}のときに満たされるので，1回のプリフェッチは
\SI{64}{\kibi\byte} から \SI{1}{\mebi\byte} までの範囲になる．}

\begin{example}
  \label{ex:l07-tbnp}
  \SI{2}{\mebi\byte} 領域の葉 0，1，2，4 にこの順でファーフォールトが
  起きたとして，Ganguly らの規則を適用する（\cref{fig:l07-tbnp-tree}）．
  各フォールトの後，GMMU はすべての祖先の有効サイズを合計する．半分を
  超えて有効なノードは，空いているほうの子を満ちているほうの子の
  サイズまで満たす．
  \begin{center}
    \small
    \begin{tabular}{@{}c>{\raggedright\arraybackslash}p{58mm}c@{}}
      \toprule
      フォールト & 50\,\% を超える祖先 & 移動量 \\
      \midrule
      葉 0 & なし（128-KiB ノード：128 のうち 64） & \SI{64}{\kibi\byte} \\
      葉 1 & 128-KiB ノードは満杯で満たすものなし（256-KiB ノード：256 のうち 128） & \SI{64}{\kibi\byte} \\
      葉 2 & 256-KiB ノード：256 のうち 192．葉 3 を満たす & \SI{128}{\kibi\byte} \\
      葉 4 & 512-KiB ノード：512 のうち 320．葉 5--7 を満たす & \SI{256}{\kibi\byte} \\
      \bottomrule
    \end{tabular}
  \end{center}
  4回のフォールトで，4回のオンデマンドな \SI{4}{\kibi\byte} の移動に
  よる \SI{16}{\kibi\byte} ではなく，\SI{512}{\kibi\byte}，8枚の葉が
  持ち込まれる．葉0だけでは自分のノードを越えるプリフェッチは
  起こらなかった．ちょうど 50\,\% ではノードは満たされない．
\end{example}

\begin{figure}[tb]
  \centering
  % Tree-based neighbourhood prefetching (Lesson 7): a 2 MiB region of a
% managed allocation as a full binary tree whose 32 leaves are 64 KiB
% basic blocks (16 pages of 4 KiB).  Shaded leaves show the end state
% of the worked example in the text: far faults in leaves 0, 1, 2 and
% 4 (red) and the leaves prefetched by the >50 % rule (green).
% Redraws the outline figure m7l2_treebase, with the intermediate
% levels of Ganguly et al. (2019) drawn in.
% Usage: % Tree-based neighbourhood prefetching (Lesson 7): a 2 MiB region of a
% managed allocation as a full binary tree whose 32 leaves are 64 KiB
% basic blocks (16 pages of 4 KiB).  Shaded leaves show the end state
% of the worked example in the text: far faults in leaves 0, 1, 2 and
% 4 (red) and the leaves prefetched by the >50 % rule (green).
% Redraws the outline figure m7l2_treebase, with the intermediate
% levels of Ganguly et al. (2019) drawn in.
% Usage: % Tree-based neighbourhood prefetching (Lesson 7): a 2 MiB region of a
% managed allocation as a full binary tree whose 32 leaves are 64 KiB
% basic blocks (16 pages of 4 KiB).  Shaded leaves show the end state
% of the worked example in the text: far faults in leaves 0, 1, 2 and
% 4 (red) and the leaves prefetched by the >50 % rule (green).
% Redraws the outline figure m7l2_treebase, with the intermediate
% levels of Ganguly et al. (2019) drawn in.
% Usage: \input{figures/l07-tbnp-tree}
% Width: body (120 mm); the drawing is about 117 mm wide.
\begin{tikzpicture}[x=1mm, y=1mm, font=\scriptsize\sffamily]
  \def\w{3.3}
  % level labels
  \foreach \h/\lab in {0/64\,KiB, 1/128\,KiB, 2/256\,KiB, 3/512\,KiB, 4/1\,MiB, 5/2\,MiB} {
    \node[anchor=east, text=ln-muted] at (11,\h*6+2) {\lab};
  }
  % internal nodes and edges
  \foreach \h in {1,...,5} {
    \pgfmathtruncatemacro{\n}{2^(5-\h)-1}
    \pgfmathsetmacro{\span}{2^\h*\w}
    \foreach \j in {0,...,\n} {
      \pgfmathsetmacro{\x}{12+\j*\span+\span/2}
      \pgfmathsetmacro{\xl}{12+\j*\span+\span/4}
      \pgfmathsetmacro{\xr}{12+\j*\span+3*\span/4}
      \ifnum\h=1
        \draw[ln-rule] (\x,\h*6+2) -- (\xl,4.5);
        \draw[ln-rule] (\x,\h*6+2) -- (\xr,4.5);
      \else
        \draw[ln-rule] (\x,\h*6+2) -- (\xl,\h*6-4);
        \draw[ln-rule] (\x,\h*6+2) -- (\xr,\h*6-4);
      \fi
      \fill[ln-accent] (\x,\h*6+2) circle (0.7);
    }
  }
  % leaves
  \foreach \i in {0,...,31} {
    \def\col{white}
    \ifnum\i<3 \def\col{ln-caution!35}\fi
    \ifnum\i=4 \def\col{ln-caution!35}\fi
    \ifnum\i=3 \def\col{ln-tip!40}\fi
    \ifnum\i>4 \ifnum\i<8 \def\col{ln-tip!40}\fi\fi
    \draw[fill=\col] (12+\i*\w,-0.5) rectangle ++(\w,5);
  }
  \foreach \i in {0,4,8,12,16,20,24,28} {
    \node[anchor=north, text=ln-muted] at (12+\i*\w+\w/2,-0.6) {\i};
  }
  % legend
  \draw[fill=ln-caution!35] (80,30) rectangle ++(3,3);
  \node[anchor=west] at (83.5,31.5) {\iflangja{ファーフォールト}{far fault}};
  \draw[fill=ln-tip!40] (98,30) rectangle ++(3,3);
  \node[anchor=west] at (101.5,31.5) {\iflangja{プリフェッチ}{prefetched}};
\end{tikzpicture}

% Width: body (120 mm); the drawing is about 117 mm wide.
\begin{tikzpicture}[x=1mm, y=1mm, font=\scriptsize\sffamily]
  \def\w{3.3}
  % level labels
  \foreach \h/\lab in {0/64\,KiB, 1/128\,KiB, 2/256\,KiB, 3/512\,KiB, 4/1\,MiB, 5/2\,MiB} {
    \node[anchor=east, text=ln-muted] at (11,\h*6+2) {\lab};
  }
  % internal nodes and edges
  \foreach \h in {1,...,5} {
    \pgfmathtruncatemacro{\n}{2^(5-\h)-1}
    \pgfmathsetmacro{\span}{2^\h*\w}
    \foreach \j in {0,...,\n} {
      \pgfmathsetmacro{\x}{12+\j*\span+\span/2}
      \pgfmathsetmacro{\xl}{12+\j*\span+\span/4}
      \pgfmathsetmacro{\xr}{12+\j*\span+3*\span/4}
      \ifnum\h=1
        \draw[ln-rule] (\x,\h*6+2) -- (\xl,4.5);
        \draw[ln-rule] (\x,\h*6+2) -- (\xr,4.5);
      \else
        \draw[ln-rule] (\x,\h*6+2) -- (\xl,\h*6-4);
        \draw[ln-rule] (\x,\h*6+2) -- (\xr,\h*6-4);
      \fi
      \fill[ln-accent] (\x,\h*6+2) circle (0.7);
    }
  }
  % leaves
  \foreach \i in {0,...,31} {
    \def\col{white}
    \ifnum\i<3 \def\col{ln-caution!35}\fi
    \ifnum\i=4 \def\col{ln-caution!35}\fi
    \ifnum\i=3 \def\col{ln-tip!40}\fi
    \ifnum\i>4 \ifnum\i<8 \def\col{ln-tip!40}\fi\fi
    \draw[fill=\col] (12+\i*\w,-0.5) rectangle ++(\w,5);
  }
  \foreach \i in {0,4,8,12,16,20,24,28} {
    \node[anchor=north, text=ln-muted] at (12+\i*\w+\w/2,-0.6) {\i};
  }
  % legend
  \draw[fill=ln-caution!35] (80,30) rectangle ++(3,3);
  \node[anchor=west] at (83.5,31.5) {\iflangja{ファーフォールト}{far fault}};
  \draw[fill=ln-tip!40] (98,30) rectangle ++(3,3);
  \node[anchor=west] at (101.5,31.5) {\iflangja{プリフェッチ}{prefetched}};
\end{tikzpicture}

% Width: body (120 mm); the drawing is about 117 mm wide.
\begin{tikzpicture}[x=1mm, y=1mm, font=\scriptsize\sffamily]
  \def\w{3.3}
  % level labels
  \foreach \h/\lab in {0/64\,KiB, 1/128\,KiB, 2/256\,KiB, 3/512\,KiB, 4/1\,MiB, 5/2\,MiB} {
    \node[anchor=east, text=ln-muted] at (11,\h*6+2) {\lab};
  }
  % internal nodes and edges
  \foreach \h in {1,...,5} {
    \pgfmathtruncatemacro{\n}{2^(5-\h)-1}
    \pgfmathsetmacro{\span}{2^\h*\w}
    \foreach \j in {0,...,\n} {
      \pgfmathsetmacro{\x}{12+\j*\span+\span/2}
      \pgfmathsetmacro{\xl}{12+\j*\span+\span/4}
      \pgfmathsetmacro{\xr}{12+\j*\span+3*\span/4}
      \ifnum\h=1
        \draw[ln-rule] (\x,\h*6+2) -- (\xl,4.5);
        \draw[ln-rule] (\x,\h*6+2) -- (\xr,4.5);
      \else
        \draw[ln-rule] (\x,\h*6+2) -- (\xl,\h*6-4);
        \draw[ln-rule] (\x,\h*6+2) -- (\xr,\h*6-4);
      \fi
      \fill[ln-accent] (\x,\h*6+2) circle (0.7);
    }
  }
  % leaves
  \foreach \i in {0,...,31} {
    \def\col{white}
    \ifnum\i<3 \def\col{ln-caution!35}\fi
    \ifnum\i=4 \def\col{ln-caution!35}\fi
    \ifnum\i=3 \def\col{ln-tip!40}\fi
    \ifnum\i>4 \ifnum\i<8 \def\col{ln-tip!40}\fi\fi
    \draw[fill=\col] (12+\i*\w,-0.5) rectangle ++(\w,5);
  }
  \foreach \i in {0,4,8,12,16,20,24,28} {
    \node[anchor=north, text=ln-muted] at (12+\i*\w+\w/2,-0.6) {\i};
  }
  % legend
  \draw[fill=ln-caution!35] (80,30) rectangle ++(3,3);
  \node[anchor=west] at (83.5,31.5) {\iflangja{ファーフォールト}{far fault}};
  \draw[fill=ln-tip!40] (98,30) rectangle ++(3,3);
  \node[anchor=west] at (101.5,31.5) {\iflangja{プリフェッチ}{prefetched}};
\end{tikzpicture}

  \caption{\SI{2}{\mebi\byte} 領域のプリフェッチツリー：
    \SI{64}{\kibi\byte}（各16ページ）の32枚の葉が，ルートまでの二分の
    段の下にある．網掛けは\cref{ex:l07-tbnp}の最終状態であり，葉 0，1，
    2，4 のファーフォールトとプリフェッチされた葉を示す．}
  \label{fig:l07-tbnp-tree}
\end{figure}

\subsection{オーバーサブスクリプションと追い出し}

\begin{definition}[オーバーサブスクリプションと追い出し]
  \label{def:l07-oversubscription}
  GPU の物理メモリが保持できるより多くのマネージドメモリを割り当てる
  ことを
  \term[めもりのおーばーさぶすくりぷしょん]{メモリのオーバーサブスクリプション}[memory oversubscription]
  という．このとき場所を空けるためにページがデバイスメモリから出て
  いかなければならず，これが
  \term[ぺーじおいだし]{ページ追い出し}[page eviction]である．
  どのページが出ていくかはポリシが決める．
\end{definition}

\margin{CPG §19.1.2.7 GPU Memory Oversubscription：ユニファイドメモリでは，
プログラムが\enquote{allocate and share arrays larger than the
memory capacity of any individual processor in the system}ことが
できる \cite{nvidia2024cudaguide}．}%
プリフェッチと追い出しは相互に作用する．積極的なプリフェッチはまだ
必要なページを追い出し，プリフェッチャを無視した追い出しポリシは
散在する \SI{4}{\kibi\byte} ページを追い出して，プリフェッチャが頼る
連続した領域を壊す．この問題はカーネルのメモリ需要とともに大きくなる．
1つの答えは\term[じぜんおいだし]{事前追い出し}[pre-eviction]である．
LRU リストが追い出しに選んだページとともに，その周囲の連続した
ページを，プリフェッチャがそれらを持ち込んだ \SI{64}{\kibi\byte}
ブロックの倍数単位で追い出し，プリフェッチャが再び使える単位で空き
領域が生まれるようにする．\margin{Ganguly ら：ツリーベースの事前
追い出しをハードウェアプリフェッチャと組み合わせると，
\SI{4}{\kibi\byte} および \SI{2}{\mebi\byte} ページの LRU 置換に対して
平均 93\,\% と 18.5\,\% の高速化が得られる
\cite{ganguly2019interplay}．}

\begin{keypoint}
  共有アドレス空間は明示的なコピーを取り除くが，データの移動は取り
  除かない．ディスクリート GPU では，ページへの最初のアクセスはすべて
  ファーフォールトとなり，PCIe を越えてホストへ行って戻り，数十マイクロ秒
  かかる．各フォールトの周囲のプリフェッチ（ツリーベースの方式）は，
  スレッドブロックが順序どおりに実行されなくても往復のコストを償却する．
  オーバーサブスクリプションのもとでは，追い出しポリシはプリフェッチャが
  扱う単位を尊重しなければならない．
\end{keypoint}

\section{ワープスケジューリング}
\label{sec:l07-scheduling}

第4章のパイプラインの説明は，一度に準備完了のワープが1つであると
仮定していた．\source{M7，GTO とラウンドロビン；A18 ch.~3；K12 §1.4.1
GPU Pipeline；CPG §4.2 Hardware Multithreading；\cite{rogers2012ccws,
narasiman2011twolevel,gebhart2011energy}．}通常は複数のワープが同時に
スコアボード上ですべてのオペランドをそろえており，その中から1つを
選ばなければならない．本節では基本的なポリシを定義し，その選択が
レイテンシ隠蔽と L1 キャッシュにどう影響するかを示し，これらの観察の
上に築かれた研究上のスケジューラを紹介する．

\begin{definition}[ワープスケジューラ]
  \label{def:l07-warp-scheduler}
  \term[わーぷすけじゅーら]{ワープスケジューラ}[warp scheduler]は，
  毎サイクル，次の命令が準備完了のワープの中から1つを選び，その命令を
  発行するパイプライン段である．
\end{definition}

\margin{CPG §16.6.1（Compute Capability 7.x，Architecture）：SM は
4つのワープスケジューラを持ち，ワープをそれらに静的に分配する．各
スケジューラは準備完了のワープの1つについて毎サイクル1命令を発行する
\cite{nvidia2024cudaguide}．CC 7.0 の常駐ワープ 64 個（CPG Table 21；
CC 7.5 では 32 個）では，スケジューラあたり最大 16 ワープになる
（本書の計算）．}%
\subsection{命令スケジューリングからワープスケジューリングへ}

アウトオブオーダの CPU にも似たユニットがある．その命令スケジューラは
スケジューリングキューの準備完了の命令の中から選ぶ．既定のポリシとなるのは
\term[さいこゆうせん]{最古優先}[oldest-first]であり，準備完了の命令の
うち最初にキューに入ったものが発行される．別の方式は\emph{クリティカル}な
命令を優先する．ロードを先にしてアドレスがわかり次第メモリ要求を
始めるか，依存する命令の多い命令を先にしてより多くの命令が早く準備
完了になるようにする．

GPU のワープスケジューラは各ワープの中ではインオーダである．ワープの
命令はプログラム順に発行される．ワープをまたいでは任意の準備完了の
ワープを選べるので，SM 全体の命令の流れはアウトオブオーダに見える．
それぞれ命令 I1，I2，I3 を実行する4つのワープは，たとえば次のように
発行されうる．
\begin{center}
  \small
  I1\,w1, I1\,w2, I2\,w1, I1\,w3, I2\,w2, I3\,w1, I1\,w4, I2\,w3,
  I2\,w4, I3\,w2, I3\,w4, I3\,w3
\end{center}
これはどのワープの中でも I1，I2，I3 の順序を保っている．\margin{K12
§1.4.1：スケジューラはインオーダだが\enquote{can select any warps
that are ready}ので，SM の実行は\enquote{might look like an
out-of-order execution} \cite{kim2012gpgpu}．}どの準備完了のワープを
先にするかがポリシである．

\subsection{ラウンドロビンと GTO}

\begin{definition}[ラウンドロビンと GTO]
  \label{def:l07-rr-gto}
  \term[らうんどろびん]{ラウンドロビン}[round-robin (RR)]（RR）の
  スケジューラは，固定の巡回順でワープに順番を与え，1命令ずつ発行させる．
  \term[るーずらうんどろびん]{ルーズラウンドロビン}[loose round-robin (LRR)]（LRR）
  の形では，自分の番に発行できないワープは飛ばされ，順番で次のワープに
  機会が回る．
  \term[じーてぃーおーすけじゅーりんぐ]{GTOスケジューリング}[greedy-then-oldest (GTO)]
  （greedy-then-oldest）のスケジューラは，1つのワープが長レイテンシの
  事象（キャッシュミス，分岐，長いテンソルコア演算）でストールするまで
  そのワープから発行し続け，それから最も古い準備完了のワープに切り
  替える．ワープの年齢とは，そのワープが SM に割り当てられた時刻で
  ある．\index{RR|see{ラウンドロビン}}\index{LRR|see{ルーズラウンドロビン}}\index{GTO|see{GTOスケジューリング}}
\end{definition}

\margin{Rogers らは両方の基準方式を定義している．GTO は\enquote{runs
a single wavefront until it stalls then picks the oldest ready
wavefront}であり，年齢は\enquote{the time it is assigned to the
core}である \cite{rogers2012ccws}．}%
2つのポリシは，ワープがメモリ命令に到達する時期が異なる．
ラウンドロビンではすべてのワープがそろって進み，ほぼ同時にロードに
到達する．GTO では1つのワープが先行して早くロードを始め，他のワープが
1つずつ後に続く．

\begin{example}
  \label{ex:l07-rr-gto}
  4つのワープが同じプログラムを実行する．2つの算術命令 \texttt{c1}，
  \texttt{c2}，キャッシュでミスして 10 サイクル後に戻るロード
  \texttt{m3}，そして3つの算術命令 \texttt{c4}--\texttt{c6} であり，
  \texttt{c4} はロードした値を使う．\margin{記法：毎サイクル1命令が
発行される．サイクル $t$ に発行されたミスにより，依存する命令は
サイクル $t+10$ に発行できる．本書の数値．}\cref{fig:l07-rr-gto}は
  発行スロットを示す．
  \begin{itemize}
    \item \emph{ラウンドロビン}：サイクル 0--7 がすべてのワープの
      \texttt{c1} と \texttt{c2} を，サイクル 8--11 が4つのロードを
      発行する．すべてのワープが自分のロードを待つことになり，最初の
      ロードはサイクル 18 に戻るので，サイクル 12--17 はストールである．
      残りの 12 命令はサイクル 18--29 を要する．30 サイクル，そのうち
      6 サイクルがストールである．
    \item \emph{GTO}：ワープ1はサイクル 0--2 に \texttt{c1}，
      \texttt{c2}，\texttt{m3} を発行してストールする．ワープ2が
      サイクル 3--5，ワープ3が 6--8，ワープ4が 9--11 に続く．ワープ1の
      ロードはワープ4のロードが発行された1サイクル後のサイクル 12 に
      戻り，各ワープはデータが届くにつれて残りの3命令を終える．
      24 サイクル，ストールなしである．
  \end{itemize}
  GTO ではロードが分散し，他のワープの算術命令がそれぞれのロードを
  隠蔽する．
\end{example}

\begin{figure}[tb]
  \centering
  % Round-robin versus greedy-then-oldest issue (Lesson 7), the worked
% example of the text: four warps each run c1 c2 m3 c4 c5 c6, where m3
% is a load that misses (latency 10 cycles) and c4 needs its value.
% One instruction issues per cycle; the number in a cell is the warp,
% M marks the load, grey cells are stall cycles, and the bars below a
% row show each load's latency.
% Redraws the outline figure m7l3_rr_gto with the text's own program.
% Usage: % Round-robin versus greedy-then-oldest issue (Lesson 7), the worked
% example of the text: four warps each run c1 c2 m3 c4 c5 c6, where m3
% is a load that misses (latency 10 cycles) and c4 needs its value.
% One instruction issues per cycle; the number in a cell is the warp,
% M marks the load, grey cells are stall cycles, and the bars below a
% row show each load's latency.
% Redraws the outline figure m7l3_rr_gto with the text's own program.
% Usage: % Round-robin versus greedy-then-oldest issue (Lesson 7), the worked
% example of the text: four warps each run c1 c2 m3 c4 c5 c6, where m3
% is a load that misses (latency 10 cycles) and c4 needs its value.
% One instruction issues per cycle; the number in a cell is the warp,
% M marks the load, grey cells are stall cycles, and the bars below a
% row show each load's latency.
% Redraws the outline figure m7l3_rr_gto with the text's own program.
% Usage: \input{figures/l07-rr-gto}
% Width: body (120 mm); the drawing is about 118 mm wide.
\begin{tikzpicture}[x=3.5mm, y=1mm, font=\scriptsize\sffamily]
  \colorlet{l07d1}{ln-accent}\colorlet{l07d2}{ln-caution}%
  \colorlet{l07d3}{ln-tip}\colorlet{l07d4}{ln-source}%
  % \lsevenissue{y}{list of w/m entries; 0 = stall}
  \newcommand{\lsevenissue}[2]{%
    \foreach \w/\m [count=\c from 0] in {#2} {
      \ifnum\w=0
        \draw[fill=ln-rule!40, draw=ln-rule] (\c+3,#1) rectangle ++(1,5);
      \else
        \draw[fill=l07d\w!22] (\c+3,#1) rectangle ++(1,5)
          node[midway] {\ifnum\m=1 M\else\w\fi};
      \fi
    }}
  % \lsevenlat{y}{issue cycle}{warp}
  \newcommand{\lsevenlat}[3]{%
    \draw[|-|, thick, draw=l07d#3] (#2+3.5,#1) -- (#2+13.5,#1);}
  % round-robin
  \iflangja{\node[anchor=east, align=right, inner sep=0.3pt] at (2.8,30.5) {ラウンド\\ロビン};}%
    {\node[anchor=east, align=right] at (2.8,30.5) {round-\\robin};}
  \lsevenissue{28}{1/0,2/0,3/0,4/0,1/0,2/0,3/0,4/0,1/1,2/1,3/1,4/1,0/0,0/0,0/0,0/0,0/0,0/0,1/0,2/0,3/0,4/0,1/0,2/0,3/0,4/0,1/0,2/0,3/0,4/0}
  \lsevenlat{26}{8}{1}\lsevenlat{24.5}{9}{2}\lsevenlat{23}{10}{3}\lsevenlat{21.5}{11}{4}
  % GTO
  \node[anchor=east, align=right] at (2.8,11.5) {GTO};
  \lsevenissue{9}{1/0,1/0,1/1,2/0,2/0,2/1,3/0,3/0,3/1,4/0,4/0,4/1,1/0,1/0,1/0,2/0,2/0,2/0,3/0,3/0,3/0,4/0,4/0,4/0}
  \lsevenlat{7}{2}{1}\lsevenlat{5.5}{5}{2}\lsevenlat{4}{8}{3}\lsevenlat{2.5}{11}{4}
  % cycle axis
  \foreach \c in {0,5,10,15,20,25,30} {
    \draw[ln-muted] (\c+3,0) -- ++(0,-1) node[below, inner sep=1pt] {\c};
  }
  \draw[->, ln-muted] (3,0) -- (34,0) node[above left, inner sep=1pt] {\iflangja{サイクル}{cycle}};
\end{tikzpicture}

% Width: body (120 mm); the drawing is about 118 mm wide.
\begin{tikzpicture}[x=3.5mm, y=1mm, font=\scriptsize\sffamily]
  \colorlet{l07d1}{ln-accent}\colorlet{l07d2}{ln-caution}%
  \colorlet{l07d3}{ln-tip}\colorlet{l07d4}{ln-source}%
  % \lsevenissue{y}{list of w/m entries; 0 = stall}
  \newcommand{\lsevenissue}[2]{%
    \foreach \w/\m [count=\c from 0] in {#2} {
      \ifnum\w=0
        \draw[fill=ln-rule!40, draw=ln-rule] (\c+3,#1) rectangle ++(1,5);
      \else
        \draw[fill=l07d\w!22] (\c+3,#1) rectangle ++(1,5)
          node[midway] {\ifnum\m=1 M\else\w\fi};
      \fi
    }}
  % \lsevenlat{y}{issue cycle}{warp}
  \newcommand{\lsevenlat}[3]{%
    \draw[|-|, thick, draw=l07d#3] (#2+3.5,#1) -- (#2+13.5,#1);}
  % round-robin
  \iflangja{\node[anchor=east, align=right, inner sep=0.3pt] at (2.8,30.5) {ラウンド\\ロビン};}%
    {\node[anchor=east, align=right] at (2.8,30.5) {round-\\robin};}
  \lsevenissue{28}{1/0,2/0,3/0,4/0,1/0,2/0,3/0,4/0,1/1,2/1,3/1,4/1,0/0,0/0,0/0,0/0,0/0,0/0,1/0,2/0,3/0,4/0,1/0,2/0,3/0,4/0,1/0,2/0,3/0,4/0}
  \lsevenlat{26}{8}{1}\lsevenlat{24.5}{9}{2}\lsevenlat{23}{10}{3}\lsevenlat{21.5}{11}{4}
  % GTO
  \node[anchor=east, align=right] at (2.8,11.5) {GTO};
  \lsevenissue{9}{1/0,1/0,1/1,2/0,2/0,2/1,3/0,3/0,3/1,4/0,4/0,4/1,1/0,1/0,1/0,2/0,2/0,2/0,3/0,3/0,3/0,4/0,4/0,4/0}
  \lsevenlat{7}{2}{1}\lsevenlat{5.5}{5}{2}\lsevenlat{4}{8}{3}\lsevenlat{2.5}{11}{4}
  % cycle axis
  \foreach \c in {0,5,10,15,20,25,30} {
    \draw[ln-muted] (\c+3,0) -- ++(0,-1) node[below, inner sep=1pt] {\c};
  }
  \draw[->, ln-muted] (3,0) -- (34,0) node[above left, inner sep=1pt] {\iflangja{サイクル}{cycle}};
\end{tikzpicture}

% Width: body (120 mm); the drawing is about 118 mm wide.
\begin{tikzpicture}[x=3.5mm, y=1mm, font=\scriptsize\sffamily]
  \colorlet{l07d1}{ln-accent}\colorlet{l07d2}{ln-caution}%
  \colorlet{l07d3}{ln-tip}\colorlet{l07d4}{ln-source}%
  % \lsevenissue{y}{list of w/m entries; 0 = stall}
  \newcommand{\lsevenissue}[2]{%
    \foreach \w/\m [count=\c from 0] in {#2} {
      \ifnum\w=0
        \draw[fill=ln-rule!40, draw=ln-rule] (\c+3,#1) rectangle ++(1,5);
      \else
        \draw[fill=l07d\w!22] (\c+3,#1) rectangle ++(1,5)
          node[midway] {\ifnum\m=1 M\else\w\fi};
      \fi
    }}
  % \lsevenlat{y}{issue cycle}{warp}
  \newcommand{\lsevenlat}[3]{%
    \draw[|-|, thick, draw=l07d#3] (#2+3.5,#1) -- (#2+13.5,#1);}
  % round-robin
  \iflangja{\node[anchor=east, align=right, inner sep=0.3pt] at (2.8,30.5) {ラウンド\\ロビン};}%
    {\node[anchor=east, align=right] at (2.8,30.5) {round-\\robin};}
  \lsevenissue{28}{1/0,2/0,3/0,4/0,1/0,2/0,3/0,4/0,1/1,2/1,3/1,4/1,0/0,0/0,0/0,0/0,0/0,0/0,1/0,2/0,3/0,4/0,1/0,2/0,3/0,4/0,1/0,2/0,3/0,4/0}
  \lsevenlat{26}{8}{1}\lsevenlat{24.5}{9}{2}\lsevenlat{23}{10}{3}\lsevenlat{21.5}{11}{4}
  % GTO
  \node[anchor=east, align=right] at (2.8,11.5) {GTO};
  \lsevenissue{9}{1/0,1/0,1/1,2/0,2/0,2/1,3/0,3/0,3/1,4/0,4/0,4/1,1/0,1/0,1/0,2/0,2/0,2/0,3/0,3/0,3/0,4/0,4/0,4/0}
  \lsevenlat{7}{2}{1}\lsevenlat{5.5}{5}{2}\lsevenlat{4}{8}{3}\lsevenlat{2.5}{11}{4}
  % cycle axis
  \foreach \c in {0,5,10,15,20,25,30} {
    \draw[ln-muted] (\c+3,0) -- ++(0,-1) node[below, inner sep=1pt] {\c};
  }
  \draw[->, ln-muted] (3,0) -- (34,0) node[above left, inner sep=1pt] {\iflangja{サイクル}{cycle}};
\end{tikzpicture}

  \caption{\cref{ex:l07-rr-gto}の発行スロット：数字はワープ，M はその
    ロード，灰色はストールサイクルである．各行の下のバーは 10 サイクルの
    ロードレイテンシである．ラウンドロビンは4つのロードを続けて発行して
    ストールし，GTO はそれらをずらす．}
  \label{fig:l07-rr-gto}
\end{figure}

GTO が常に優れているわけではない．利得はプログラムに依存する．例では，
利得はロードの後の算術命令から来ている．すべてのワープがロードの直後に
同じバリアでストールするプログラムは何も得ない．ロードレイテンシが他の
すべてのワープの算術命令の合計より長ければ GTO もストールする．
ワープがデータを共有する場合は，ラウンドロビンでそろって発行するほうが，
1つのワープのミスで他のワープの要求にも応えられることがある．

\subsection{2レベルスケジューリング}

毎サイクルすべての準備完了のワープを探索するスケジューラは，1つの
発行スロットのために数十の候補を比較し，この探索は電力を消費する．

\begin{marginfigure}
  % Two-level warp scheduling (Lesson 7): the scheduler searches only the
% small active set; a warp that meets a long-latency operation moves to
% the pending set, and a pending warp whose operands are ready refills
% the active set.
% Redraws the outline figure m7l3_two_lvl (right half).
% Usage: % Two-level warp scheduling (Lesson 7): the scheduler searches only the
% small active set; a warp that meets a long-latency operation moves to
% the pending set, and a pending warp whose operands are ready refills
% the active set.
% Redraws the outline figure m7l3_two_lvl (right half).
% Usage: % Two-level warp scheduling (Lesson 7): the scheduler searches only the
% small active set; a warp that meets a long-latency operation moves to
% the pending set, and a pending warp whose operands are ready refills
% the active set.
% Redraws the outline figure m7l3_two_lvl (right half).
% Usage: \input{figures/l07-two-level}  (in a marginfigure)
% Width: margin (50 mm); the drawing is about 46 mm wide.
\begin{tikzpicture}[x=1mm, y=1mm, font=\scriptsize\sffamily,
    >={Stealth[length=1.6mm]},
    wp/.style={draw, minimum width=8mm, minimum height=4.5mm, inner sep=1pt}]
  \node[anchor=south, text=ln-muted] at (8,38) {\iflangja{待機集合}{pending}};
  \foreach \w [count=\k from 0] in {4,5,6,7,8} {
    \node[wp, fill=ln-box] at (8,35-\k*5.5) {W\w};
  }
  \node[anchor=south, text=ln-muted] at (34,38) {\iflangja{アクティブ集合}{active}};
  \foreach \w [count=\k from 0] in {1,2,3} {
    \node[wp, fill=ln-accent!20] at (34,35-\k*5.5) {W\w};
  }
  \draw[->, ln-caution] (29.5,30) to[bend right=20] node[above, text=ln-caution, align=center] {\iflangja{長レイテンシ}{long}\\[-2pt]\iflangja{}{latency}} (12.5,30);
  \draw[->, ln-tip] (12.5,19) to[bend right=20] node[below, text=ln-tip] {\iflangja{準備完了}{ready}} (29.5,19);
  \node[draw, fill=ln-tip!25, minimum width=16mm, minimum height=5mm] (s) at (34,6) {\iflangja{スケジューラ}{scheduler}};
  \draw[->] (34,21.7) -- (s);
\end{tikzpicture}
  (in a marginfigure)
% Width: margin (50 mm); the drawing is about 46 mm wide.
\begin{tikzpicture}[x=1mm, y=1mm, font=\scriptsize\sffamily,
    >={Stealth[length=1.6mm]},
    wp/.style={draw, minimum width=8mm, minimum height=4.5mm, inner sep=1pt}]
  \node[anchor=south, text=ln-muted] at (8,38) {\iflangja{待機集合}{pending}};
  \foreach \w [count=\k from 0] in {4,5,6,7,8} {
    \node[wp, fill=ln-box] at (8,35-\k*5.5) {W\w};
  }
  \node[anchor=south, text=ln-muted] at (34,38) {\iflangja{アクティブ集合}{active}};
  \foreach \w [count=\k from 0] in {1,2,3} {
    \node[wp, fill=ln-accent!20] at (34,35-\k*5.5) {W\w};
  }
  \draw[->, ln-caution] (29.5,30) to[bend right=20] node[above, text=ln-caution, align=center] {\iflangja{長レイテンシ}{long}\\[-2pt]\iflangja{}{latency}} (12.5,30);
  \draw[->, ln-tip] (12.5,19) to[bend right=20] node[below, text=ln-tip] {\iflangja{準備完了}{ready}} (29.5,19);
  \node[draw, fill=ln-tip!25, minimum width=16mm, minimum height=5mm] (s) at (34,6) {\iflangja{スケジューラ}{scheduler}};
  \draw[->] (34,21.7) -- (s);
\end{tikzpicture}
  (in a marginfigure)
% Width: margin (50 mm); the drawing is about 46 mm wide.
\begin{tikzpicture}[x=1mm, y=1mm, font=\scriptsize\sffamily,
    >={Stealth[length=1.6mm]},
    wp/.style={draw, minimum width=8mm, minimum height=4.5mm, inner sep=1pt}]
  \node[anchor=south, text=ln-muted] at (8,38) {\iflangja{待機集合}{pending}};
  \foreach \w [count=\k from 0] in {4,5,6,7,8} {
    \node[wp, fill=ln-box] at (8,35-\k*5.5) {W\w};
  }
  \node[anchor=south, text=ln-muted] at (34,38) {\iflangja{アクティブ集合}{active}};
  \foreach \w [count=\k from 0] in {1,2,3} {
    \node[wp, fill=ln-accent!20] at (34,35-\k*5.5) {W\w};
  }
  \draw[->, ln-caution] (29.5,30) to[bend right=20] node[above, text=ln-caution, align=center] {\iflangja{長レイテンシ}{long}\\[-2pt]\iflangja{}{latency}} (12.5,30);
  \draw[->, ln-tip] (12.5,19) to[bend right=20] node[below, text=ln-tip] {\iflangja{準備完了}{ready}} (29.5,19);
  \node[draw, fill=ln-tip!25, minimum width=16mm, minimum height=5mm] (s) at (34,6) {\iflangja{スケジューラ}{scheduler}};
  \draw[->] (34,21.7) -- (s);
\end{tikzpicture}

  \caption{2レベルスケジューリング：スケジューラはアクティブ集合だけを
    探索する．}
  \label{fig:l07-two-level}
\end{marginfigure}

\begin{definition}[2レベルワープスケジューリング]
  \label{def:l07-two-level}
  \term[にれべるわーぷすけじゅーりんぐ]{2レベルワープスケジューリング}[two-level warp scheduling]
  では，ワープは
  \term[あくてぃぶわーぷしゅうごうたいきわーぷしゅうごう]{アクティブワープ集合／待機ワープ集合}[active warp set / pending warp set]
  の組に分けられる．探索されるのは少数のアクティブなワープだけである．
  長レイテンシの演算に出会ったワープは待機集合に移り，演算が完了した
  待機中のワープは，空きがあればアクティブ集合に戻る
  （\cref{fig:l07-two-level}）．
\end{definition}

この考えには2つの論文が引かれ，その動機は異なる．\caution{講義 vs
論文：Gebhart ら \cite{gebhart2011energy} はアクティブ／待機の分割を
\emph{エネルギー}のために導入している（レジスタファイルキャッシュと
合わせて，レジスタファイルのエネルギーを 36\,\% 削減）．Narasiman ら
\cite{narasiman2011twolevel} は\emph{性能}のためにワープをフェッチ
グループに分ける．ラージワープと合わせて 19.1\,\% である．}Narasiman
らはワープを\term[ふぇっちぐるーぷ]{フェッチグループ}[fetch group]に
分ける．たとえば 32 ワープを8ワープずつの4グループにする．1つの
グループは，そのすべてのワープがストールするまで優先権を持ち，グループ
内でもグループ間でもラウンドロビンである．こうしてグループはそれぞれ
異なる時期に長レイテンシの演算に到達する．これは\cref{ex:l07-rr-gto}と
同じずらしであり，あるグループがメモリを待つ間，常にどこかのグループに
仕事がある．

\subsection{スケジューリングとキャッシュの振る舞い}

スケジューラは，どのデータが同時に L1 キャッシュにあるべきかも決める．

\begin{definition}[ワーキングセット，スラッシング，初期参照ミス]
  \label{def:l07-working-set}
  ある区間におけるプログラムの
  \term[わーきんぐせっと]{ワーキングセット}[working set]とは，その区間に
  プログラムが使うメモリブロックの集合である．ワーキングセットが
  キャッシュを超え，ブロックが再利用される前に追い出されてほとんどの
  アクセスがミスするとき，キャッシュは
  \term[きゃっしゅすらっしんぐ]{キャッシュスラッシング}[cache thrashing]
  を起こしている．\term[しょきさんしょうみす]{初期参照ミス}[compulsory miss]
  はブロックへの最初のアクセスのミスであり，どのキャッシュも避けられない．
\end{definition}

同じワープがキャッシュに持ち込んだブロックへのヒットは
\term[わーぷないきょくしょせい]{ワープ内局所性}[intra-warp locality]を，
別のワープが持ち込んだブロックへのヒットは
\term[わーぷかんきょくしょせい]{ワープ間局所性}[inter-warp locality]を
利用している．
\margin{Rogers らはベンチマークのヒットをこのように分類している．
キャッシュに強く依存するものでは，再利用のほとんどがワープ内である
\cite{rogers2012ccws}．}次の例のように各ワープが自分のデータを再利用する
とき，ラウンドロビンは最悪の順序である．

\begin{example}
  \label{ex:l07-thrashing}
  3つのワープがそれぞれ自分のブロックを4回読む．ワープ $k$ はブロック
  B$k$ を読む．L1 は LRU 置換（第6章）で2ブロックを保持し，各アクセスは
  次のアクセスが発行される前にヒットかミスかが決まる．
  \begin{center}
    \small
    \begin{tabular}{@{}ll@{}}
      \toprule
      順序 & アクセス（太字＝ミス） \\
      \midrule
      RR  & \textbf{B1 B2 B3 B1 B2 B3 B1 B2 B3 B1 B2 B3} \\
      GTO & \textbf{B1} B1 B1 B1 \textbf{B2} B2 B2 B2 \textbf{B3} B3 B3 B3 \\
      \bottomrule
    \end{tabular}
  \end{center}
  ラウンドロビンでは，各ブロックはそのワープが戻ってくる前に他の2つに
  追い出される．12 回のアクセスで 12 回のミスである．GTO では各ワープが
  次のワープの開始前に自分のブロックを終えるので，3回の初期参照ミス
  だけである．\caution{GTO の行は，ワープが最初のミスの後も発行を
続けると仮定している．そのミスの時点で他のワープが準備完了なら GTO は
切り替え，そのワーキングセットは複数のワープのデータを含む．そのため
下の式で GTO の $\mathit{SW} \ge 1$ となる．}
\end{example}

一般に，各ワープのワーキングセットの大きさを $\mathit{SS}_W$ とし，
SM に $N$ 個のワープが常駐するとする．ラウンドロビンではすべての $N$
ワープが同時に進行中であり，L1 は $\mathit{SS}_W \times N$ を保持
する必要がある．スケジュール可能なワープ数を $\mathit{SW}$ に
制限すれば $\mathit{SS}_W \times \mathit{SW}$ で済む．
GTO は長レイテンシの事象のたびにワープを切り替えるので，
$\mathit{SW}$ は必ずしも1ではない．\cref{ex:l07-thrashing}では
$\mathit{SS}_W = 1$ ブロック，$N = 3$ で，2ブロックのキャッシュを
超える．$\mathit{SW} = 2$ ならワーキングセットは収まる．\tip{必要な
L1 $\approx$ $\mathit{SS}_W$ と発行しうるワープ数の積：ラウンドロビン
では $N$，制限下では $\mathit{SW}$．大きさ $C$ のキャッシュに収める
には $\mathit{SW} \le C/\mathit{SS}_W$ に制限する．}

\begin{definition}[キャッシュ考慮型ウェーブフロントスケジューリング]
  \label{def:l07-ccws}
  Rogers らのスケジューラは
  \term[きゃっしゅこうりょがたうぇーぶふろんとすけじゅーりんぐ]{キャッシュ考慮型ウェーブフロントスケジューリング}[cache-conscious wavefront scheduling (CCWS)]（CCWS）
  と呼ばれ \cite{rogers2012ccws}，発行できるワープの数を制限して，
  それらのワーキングセットが L1 に収まるようにする．ストールの際には
  少数のワープ（たとえば W1--W3）の中から選び，それらが完了したときに
  だけ別のワープ（W4）を開始する．\index{CCWS|see{キャッシュ考慮型ウェーブフロントスケジューリング}}
\end{definition}

\caution{講義 vs 論文：ここで述べた固定の制限は，論文の基準方式
\enquote{static wavefront limiting}である．CCWS そのものは適応的で
ある．検出器が，追い出しによってワープの局所性が壊されたことに気づくと，
スケジューラは他のワープを絞ってそのワープに L1D へのより排他的な
アクセスを与える．}%
論文の仕組みは
\term[うぇーぶふろんとないきょくしょせいそんしつけんしゅつき]{ウェーブフロント内局所性損失検出器}[lost intra-wavefront locality detector (LLD)]（LLD）
であり，ワープごとに，そのワープが L1 に持ち込み追い出しによって
失ったラインのタグ（データを持たない犠牲タグ配列）を保持する．ワープ
自身の犠牲タグにヒットする L1 ミスは失われたワープ内局所性であり，
そのワープの局所性損失スコアを引き上げる．他のワープの高いスコアに
よってスコアスタックのカットオフを越えて押し出されたワープ（論文では
カットオフより上）は，スコアが減衰するまで L1D にロードを発行できない．
報告された結果は，従来のスケジューリングポリシに対し，L1 データ
キャッシュミスが平均 25\,\% 減，調和平均 24\,\% の高速化で
ある．\margin{CCWS のハードウェアはチップ全体の面積に 0.17\,\% を
加える．単純な貪欲スケジューラに対しては，キャッシュに強く依存する
ベンチマークでの調和平均の利得は 63\,\% である \cite{rogers2012ccws}．}\index{LLD|see{ウェーブフロント内局所性損失検出器}}

ワープを制限すると2つの問題が生じる．
\begin{itemize}
  \item \emph{同期．}  バリア（第3章）は，ブロックのワープがすべて
    到着したときにだけそれらを解放する．他のワープが
    \texttt{\_\_syncthreads()} で待っているブロックのワープを
    スケジューラが留め置くと，そのブロックのどのワープも進めない．
    デッドロックである．したがって制限は，他のワープが待っている
    ワープを決して締め出してはならない．
  \item \emph{制限の大きさ．}  ワープが多いほどレイテンシを隠蔽する
    スレッドレベル並列性が増え，ワープが少ないほどワーキングセットが
    小さくなる．最良の値はキャッシュとプログラムで決まり，
    CCWS は実行時に適応させる．
\end{itemize}
\Needspace{8\baselineskip}
次の形のカーネルでは，どのワープもフェーズ2を始める前に，ブロックの
すべてのワープがバリアに到達しなければならない．

\begin{lstlisting}[language=C++, numbers=none]
__global__ void twoPhase(float *a) {
  // フェーズ1：ブロックのすべてのワープが働く
  __syncthreads();   // すべてのワープを待つ
  // フェーズ2
}
\end{lstlisting}

\subsection{その他のスケジューリングポリシ}

スケジューラは L1 より先のメモリシステムにも役立てられる．\margin{GPUMech
（第8章）は RR と GTO の両方の性能を解析的にモデル化し，平均誤差は
13.2\,\% と 14.0\,\% である \cite{huang2014gpumech}．}第4章のメモリ
チャネルと同様に，各 DRAM 内のバンクは並列に要求に応えられる．同時に
稼働しているバンクの数がその
\term[ばんくれべるへいれつせい]{バンクレベル並列性}[bank-level parallelism]
である．
\begin{itemize}
  \item \emph{プリフェッチを考慮したスケジューリング}
    \cite{jog2013orchestrated} は2レベルスケジューラを基礎とするが，
    そのグループを\emph{連続しない}ワープから作る．連続するワープは
    隣接するデータにアクセスするので，1つのワープが起動した単純な
    プリフェッチャが次のワープのデータを取ってこられる．離して
    スケジュールすれば，次のワープはプリフェッチが間に合うだけ遅れて
    実行され，同時に進行中の要求はより多くの DRAM バンクに分散する．\margin{Jog
らは，シミュレートした 30 コアの GPU で，それぞれプリフェッチ付きの
ラウンドロビンと2レベルスケジューリングに対して平均 25\,\% と 7\,\% の
向上を報告している \cite{jog2013orchestrated}．}
  \item \term[しーてぃーえーこうりょすけじゅーりんぐ]{CTA考慮スケジューリング}[CTA-aware scheduling]
    （OWL \cite{jog2013owl}）は，どの CTA（スレッドブロック，第3章）を
    一緒にスケジュールするかを決め，少数の CTA のワープをまとめて，
    そのデータがキャッシュに残るようにする．これを局所性とバンクレベル
    並列性を考慮したワープ選択と組み合わせる．\margin{OWL：シミュレート
した 28 コアの GPU で，ラウンドロビンのワープスケジューリングに対して
平均 33\,\% の向上．その動機として，38 のアプリケーションのうち 19 が
平均してサイクルの 62\,\% をどのワープも実行していない状態で費やす
\cite{jog2013owl}．}
\end{itemize}

\begin{keypoint}[章のまとめ]
  GPU の仮想メモリはスレッドではなくワープを変換する．変換の数は
  コアレッシングが決め，コアレスされないロードはスレッドごとに1回の変換を
  必要としうる．ページウォークキャッシュを持つ共有の高帯域幅ページ
  テーブルウォーカ，多段 TLB，セグメントマッピング，ラージページが
  コストを抑える．ラージページは TLB のリーチを内部断片化と交換する．
  共有仮想アドレス空間はコピー後実行をオンデマンドページングに置き
  換え，そのファーフォールトは PCIe 越しに数十マイクロ秒かかる．
  ツリーベース近傍プリフェッチがそれを償却し，オーバーサブスクリプション
  下の追い出しはプリフェッチャの単位を尊重しなければならない．
  ワープスケジューラはワープの中ではインオーダに発行するが，準備完了の
  ワープの中から選ぶ．ラウンドロビンはすべてのワープをそろって進めて
  ミスを集中させ，GTO はそれらをずらす．2レベルスケジューラは
  エネルギーと性能のために小さなアクティブ集合を探索する．そして
  スケジュール可能なワープを制限すること（CCWS）は，バリアを尊重する
  限り，L1 のワーキングセットを縮める．
\end{keypoint}

\end{document}
